\documentclass[sigplan,twocolumn,review,anonymous]{acmart}

% EuroSys 2026 recommends the official ACM SIGPLAN/acmart proceedings setup.
% Replace the submission ID and conference metadata with the final values supplied by EuroSys/ACM.
\acmSubmissionID{PAPER-ID}
\setcopyright{none}
\settopmatter{printfolios=true,printacmref=false,printccs=false}
\renewcommand\footnotetextcopyrightpermission[1]{}
\pagestyle{plain}

\acmConference[EuroSys '26]{The European Conference on Computer Systems}{April 27--30, 2026}{Edinburgh, United Kingdom}
\acmYear{2026}
\copyrightyear{2026}
\acmDOI{}
\acmISBN{}

% acmart already loads compatible math fonts; avoid reloading amssymb under newtxmath.
\usepackage{amsmath}
\usepackage{algorithm}
\usepackage{algpseudocode}
\usepackage{mathtools}
\usepackage{adjustbox}
\usepackage{xspace}
\usepackage{balance}

\allowdisplaybreaks
\emergencystretch=2em
\sloppy

\makeatletter
\newenvironment{breakablealgorithm}
{%
  \par\addvspace{\baselineskip}%
  \begingroup\footnotesize
  \setlength{\baselineskip}{10pt}%
  \refstepcounter{algorithm}%
  \noindent\hrule\kern2pt%
  \renewcommand{\caption}[1]{%
    \noindent\textbf{Algorithm~\thealgorithm} ##1\par\kern2pt\hrule\kern2pt%
  }%
}
{%
  \kern2pt\hrule\par\endgroup\addvspace{\baselineskip}%
}
\makeatother
\algrenewcommand\algorithmicindent{0.7em}

\newcommand{\func}[1]{\operatorname{\mathsf{#1}}}
\newcommand{\hash}{\func{hash}}
\newcommand{\Cert}{\func{Cert}}
\newcommand{\NotarVote}{\func{NotarVote}}
\newcommand{\FirstVote}{\func{FirstVote}}
\newcommand{\FinalVote}{\func{FinalVote}}
\newcommand{\FinalSigners}{\func{FinalSigners}}
\newcommand{\NotarCert}{\func{NotarCert}}
\newcommand{\FirstCert}{\func{FirstCert}}
\newcommand{\FinalCert}{\func{FinalCert}}
\newcommand{\TimeoutCert}{\func{TimeoutCert}}
\newcommand{\ConflictProof}{\func{ConflictProof}}
\newcommand{\manyVotes}{\func{manyVotes}}
\newcommand{\allVotes}{\func{allVotes}}
\newcommand{\maxVotes}{\func{maxVotes}}
\newcommand{\id}{\mathsf{id}}
\newcommand{\Report}{\func{Report}}
\newcommand{\Slot}{\func{Slot}}
\newcommand{\NextSlot}{\func{NextSlot}}
\newcommand{\SegmentBlocks}{\func{SegmentBlocks}}
\newcommand{\SegmentValid}{\func{SegmentValid}}
\newcommand{\VoteBatch}{\func{VoteBatch}}
\newcommand{\VoteLeaf}{\func{VoteLeaf}}
\newcommand{\ReportCommittee}{\func{ReportCommittee}}
\newcommand{\SignerCommittees}{\func{SignerCommittees}}
\newcommand{\EffectiveCommittees}{\func{EffectiveCommittees}}
\newcommand{\Close}{\func{Close}}
\newcommand{\CloseCut}{\func{CloseCut}}
\newcommand{\CloseRecord}{\func{CloseRecord}}
\newcommand{\FreshCloseHash}{\func{FreshCloseHash}}
\newcommand{\CloseDead}{\func{CloseDead}}
\newcommand{\LiveCarry}{\func{LiveCarry}}
\newcommand{\CloseConflictProof}{\func{CloseConflictProof}}
\newcommand{\closeDecision}{\mathsf{closeDecision}}
\newcommand{\closeRound}{\mathsf{closeRound}}
\newcommand{\closeFinished}{\func{closeFinished}}
\newcommand{\started}{\func{started}}
\newcommand{\firstVoted}{\func{firstVoted}}
\newcommand{\finalVoted}{\func{finalVoted}}
\newcommand{\finished}{\func{finished}}
\newcommand{\enabled}{\func{enabled}}
\newcommand{\admissible}{\func{admissible}}
\newcommand{\admissibleFresh}{\func{admissibleFresh}}
\newcommand{\admissibleHistorical}{\func{admissibleHistorical}}
\newcommand{\dom}{\func{dom}}
\newcommand{\UpdateVisible}{\func{UpdateVisible}}
\newcommand{\StartClosing}{\func{StartClosing}}
\newcommand{\FinishClosing}{\func{FinishClosing}}
\newcommand{\valueOf}{\func{value}}
\newcommand{\certOutcome}{\func{certOutcome}}
\newcommand{\Terminal}{\func{Terminal}}
\newcommand{\ExtendsTree}{\func{ExtendsTree}}
\newcommand{\CompleteBlock}{\func{CompleteBlock}}
\newcommand{\ParentComplete}{\func{ParentComplete}}
\newcommand{\Complete}{\func{Complete}}
\newcommand{\timeoutVal}{\mathsf{timeout}}
\newcommand{\pending}{\mathsf{pending}}
\newcommand{\timeoutOutcome}{\mathsf{timeout}}
\newcommand{\notarized}{\func{notarized}}
\newcommand{\converged}{\func{converged}}
\newcommand{\fast}{\func{fast}}
\newcommand{\final}{\func{final}}
\newcommand{\notarKind}{\mathsf{notar}}
\newcommand{\fastKind}{\mathsf{fast}}
\newcommand{\finalKind}{\mathsf{final}}
\newcommand{\true}{\mathsf{true}}
\newcommand{\false}{\mathsf{false}}
\newcommand{\completeTree}{\mathsf{completeTree}}
\newcommand{\blockMap}{\mathsf{block}}
\newcommand{\accountOf}{\func{account}}
\newcommand{\frontier}{\mathsf{frontier}}
\newcommand{\ledgerRoot}{\mathsf{ledgerRoot}}
\newcommand{\closeHash}{\mathsf{closeHash}}
\newcommand{\carry}{\mathsf{carry}}
\newcommand{\Frontier}{\func{Frontier}}
\newcommand{\EpochBlocks}{\func{EpochBlocks}}
\newcommand{\parent}{\func{parent}}
\newcommand{\PayloadAvailable}{\func{PayloadAvailable}}
\newcommand{\ValidBlock}{\func{ValidBlock}}
\newcommand{\slotOf}{\func{slot}}
\newcommand{\finalizedSlot}{\mathsf{finalizedSlot}}
\newcommand{\FinalizeSegment}{\func{FinalizeSegment}}
\newcommand{\committee}{\mathsf{committee}}
\newcommand{\committeeAt}{\func{committeeAt}}
\newcommand{\Committees}{\func{Committees}}
\newcommand{\deriveCommittee}{\func{DeriveCommittee}}
\newcommand{\VotesIn}{\func{VotesIn}}
\newcommand{\NotarVotesIn}{\func{NotarVotesIn}}
\newcommand{\NotarizedBy}{\func{NotarizedBy}}
\newcommand{\timeoutReady}{\func{timeoutReady}}
\newcommand{\timeoutAllowed}{\func{timeoutAllowed}}
\newcommand{\FastDead}{\func{FastDead}}
\newcommand{\FinalDead}{\func{FinalDead}}
\newcommand{\TornFor}{\func{TornFor}}
\newcommand{\AllValuesDead}{\func{AllValuesDead}}
\newcommand{\Values}{\mathcal{V}}
\newcommand{\Reportable}{\func{Reportable}}
\newcommand{\ordinaryFinalized}{\func{ordinaryFinalized}}
\newcommand{\cannotFinalize}{\func{cannotFinalize}}
\newcommand{\reportLocked}{\func{reportLocked}}
\newcommand{\recordSafe}{\func{recordSafe}}
\newcommand{\RecordSupported}{\func{RecordSupported}}
\newcommand{\LedgerIncludesMap}{\func{LedgerIncludesMap}}
\newcommand{\reportStore}{\mathsf{reportStore}}
\newcommand{\ReportEquivocated}{\func{ReportEquivocated}}
\newcommand{\EligibleReport}{\func{EligibleReport}}
\newcommand{\votedFor}{\func{votedFor}}
\newcommand{\released}{\func{released}}
\newcommand{\safeToVote}{\func{safeToVote}}
\newcommand{\activeVote}{\func{activeVote}}
\newcommand{\closeCut}{\mathsf{cut}}
\newcommand{\closeState}{\mathsf{closeState}}
\newcommand{\governingHash}{\func{governingHash}}
\newcommand{\requiredFrontier}{\func{requiredFrontier}}
\newcommand{\DescendsFrom}{\func{DescendsFrom}}
\newcommand{\safeFromClose}{\func{safeFromClose}}
\newcommand{\StablePrefix}{\func{StablePrefix}}
\newcommand{\ReportProof}{\func{ReportProof}}
\newcommand{\ReportVisible}{\func{ReportVisible}}
\newcommand{\SelectedState}{\func{SelectedState}}
\newcommand{\FinalizedState}{\func{FinalizedState}}
\newcommand{\ReleasedState}{\func{ReleasedState}}
\newcommand{\cutVersion}{\mathsf{cutVersion}}
\newcommand{\frontierVersion}{\mathsf{frontierVersion}}
\newcommand{\ReconReq}{\func{ReconReq}}
\newcommand{\ReconDelta}{\func{ReconDelta}}
\newcommand{\ReconMiss}{\func{ReconMiss}}
\newcommand{\ApplyDelta}{\func{ApplyDelta}}
\newcommand{\seqOf}{\func{sequence}}
\newcommand{\ownerKey}{\mathsf{ownerKey}}
\newcommand{\signatureOf}{\func{signature}}
\newcommand{\sendsOf}{\func{sends}}
\newcommand{\receivesOf}{\func{receives}}
\newcommand{\representativeOf}{\func{representative}}
\newcommand{\bal}{\mathsf{bal}}
\newcommand{\repOf}{\mathsf{rep}}
\newcommand{\pendingSend}{\mathsf{pendingSend}}
\newcommand{\receivedSend}{\mathsf{receivedSend}}
\newcommand{\stake}{\func{stake}}
\newcommand{\weight}{\func{weight}}
\newcommand{\SignerWeight}{\func{SignerWeight}}

\title{RAI Block-Lattice Protocol: Minimal Specification}
\author{Anonymous Authors}
\affiliation{%
  \institution{Paper under review}
  \city{}
  \country{}
}
\email{}

\begin{abstract}
RAI applies one election mechanism to account-qualified segment, close-cut, and close-record elections. A slot election certifies the tip of a contiguous owner-authorized account-chain segment and atomically finalizes every block in that segment. Logical votes may be carried in vectorized signed vote batches, while close-cut report coverage is obtained from one deterministic reporting committee and ordinary finality continues to require both election committees. This specification defines the block-lattice ledger, send and receive transitions, representative delegation, stake-weighted committees, epoch closing, safety, and liveness.
\end{abstract}

\keywords{block lattice, replication, consensus, stake-weighted elections, finality, Byzantine fault tolerance}

\begin{document}
\maketitle

\section{Overview}

RAI uses three election types.  A slot election \(\Slot(a,s,e)\) selects the tip of a contiguous
account-chain segment beginning at the account's next unfinalized slot \((a,s)\) in epoch \(e\).
The tip hash transitively commits to every intermediate block in that segment.  A close-cut election
\(\CloseCut(e,r)\) certifies a canonical set of visible segment obligations that epoch \(e\) must
drain.  A close-record election \(\CloseRecord(e,r)\) certifies the canonical finalized account
frontiers through epoch \(e\), chained to the preceding close-record hash.

Logical first, notarization, and final votes retain their election-local semantics, but a signer may
place many logical vote leaves for independent elections into one canonical Merkle-rooted signed vote
batch.  A leaf normally omits committee scope and is then interpreted identically in every applicable
committee containing the signer.  When committee-local actions differ, the leaf may instead bind an
explicit proper subset of those committees.  A one-leaf batch is the scalar encoding.  Batch transport
changes neither quorum counting nor per-committee election locks.

Only three durable authenticated object classes are disseminated: blocks, signed reports, and signed vote batches.  Certificates, report quorums, conflicts, close-round death, close cuts, frontier maps, close
records, outcomes, decisions, and installed close state are deterministic local facts derived from the
validated object store.  Correct replicas use duplicate-suppressing reliable gossip so that every valid
block, report, and vote batch learned by one correct online replica eventually becomes available to every
correct online replica.

Close votes carry only a canonical hash.  A replica retains every locally derived close-cut set and
frontier map for which it votes, indexed by that hash.  If another replica knows only a different
version, it reconstructs the target by requesting an add/remove delta from a replica that retains both
preimages, applies the delta, recomputes the target hash, and validates the reconstructed contents from
its locally available blocks, reports, and vote batches.  Reconciliation messages are transport metadata, not
consensus evidence.

A notarization never finalizes a segment by itself.  A segment becomes finalized through an ordinary
global fast/final certificate for its tip or through inclusion on an account chain committed by a
decided close-record hash.  The installation replays the segment ancestor-to-descendant in temporary
state and exposes all affected finalized-slot and ledger writes atomically.  Timeouts are local to one
election and never skip the starting account slot.  A retry of the same starting slot is permitted only
after certified release.  While an obligation is unresolved, no later slot election for that account is
enabled; after finality, the next election begins immediately after the certified tip.  Receive
references remain restricted to sends finalized before the segment, preventing atomic segments from
creating cross-account cycles or speculative funding.  Elections at distance at least two epochs are
protected because every vote in epoch \(e\) commits to and obeys the certified close state and finalized
account frontiers of epoch \(e-2\).  Safety further requires authenticated optional vote scopes with
deterministic effective-committee derivation, one first vote per correct replica and committee instance,
satisfaction of every targeted committee-local voting guard before emission, a certified segment base,
one finalized block per account slot, and certified release before a correct replica votes for a
conflicting retry.

\section{Objects, State, and Committees}

An account slot is a pair \(\lambda=(a,s)\), where \(a\in\mathsf{Accounts}\) and
\(s\ge1\) is the account-local sequence number.  An election id is either
\begin{center}
\begin{adjustbox}{max width=\columnwidth}
\(\displaystyle
    \id=\Slot(a,s,e),\qquad \id=\CloseCut(e,r),\qquad\text{or}\qquad
    \id=\CloseRecord(e,r).
\)
\end{adjustbox}
\end{center}
A slot-election non-timeout value is a segment-tip block id \(b=\hash(B_t)\), where
\(\slotOf(B_t)=(a,t)\) for some \(t\ge s\).  The candidate segment is the unique contiguous parent-chain
suffix \(B_s,\ldots,B_t\) beginning immediately after the certified finalized frontier of account
\(a\).  The election is enabled only when \(s=\NextSlot(a)\), so no separately enabled election can
overlap an unresolved segment.  Because each block commits to its predecessor, the voted tip hash,
election id, and certified base commit transitively to the complete segment.  A one-block segment
\(t=s\) is the scalar special case.

Every account \(a\) has an owner verification key \(\ownerKey[a]\), with corresponding owner signing
key \(sk_a\), fixed by genesis or by an authenticated account-creation rule.  Every non-genesis block
in a candidate segment must carry a signature verifiable under \(\ownerKey[a]\).

A block has canonical body
\[
\operatorname{body}(B)=
\bigl(\slotOf(B),\parent(B),\sendsOf(B),\receivesOf(B),\representativeOf(B)\bigr).
\]
The send list contains entries \((d,x)\), where \(d\) is a destination account and \(x\) is a positive
integer amount.  A send is uniquely referenced by \(u=(\hash(B_s),i)\), where \(i\) is its canonical
index in \(\sendsOf(B_s)\).  The receive list contains such references.  The representative field is a
registered replica id.  The owner signature is
\[
\signatureOf(B)=\operatorname{Sign}_{sk_a}
\bigl(\hash(\mathsf{RAI/AccountBlock/v1}\parallel\operatorname{canon}(\operatorname{body}(B)))\bigr),
\]
and \(\hash(B)\) commits to both the canonical body and this signature.

The ledger state derived from a finalized block set is
\[
L=(\bal,\repOf,\pendingSend,\receivedSend).
\]
Here \(\bal[a]\) is the spendable balance of account \(a\), \(\repOf[a]\) is its latest
representative, \(\pendingSend[u]=(d,x)\) records an unreceived send, and
\(\receivedSend[u]\) records a consumed send reference.  This state is derived by replay from genesis;
it is not an independently trusted object.

For epoch \(e\), a close-cut version is a canonically ordered set
\[
C\subseteq\{\Slot(a,s,e):a\in\mathsf{Accounts},\ s\ge1\}.
\]
Its close-cut value is
\[
h_C=\hash\bigl(\mathsf{RAI/CloseCut}\parallel e\parallel\operatorname{canon}(C)\bigr).
\]
The voted value is only \(h_C\); the set \(C\) is local derived state retained as
\(\cutVersion[e,h_C]=C\).

For epoch \(e\), a frontier version is the canonically ordered map
\[
\Phi=\operatorname{sort}\{(a,f(a)):a\in\mathsf{Accounts}\},
\]
where every \(f(a)\) is a block hash.  The close-record value is
\[
h_F=\hash\bigl(
\mathsf{RAI/CloseRecord}\parallel e\parallel\closeHash[e-1]
\parallel\operatorname{canon}(\Phi)
\bigr).
\]
The voted value is only \(h_F\); the map \(\Phi\) is local derived state retained as
\(\frontierVersion[e,h_F]=\Phi\).  The local ledger-state root is the convenience value
\[
R^{\mathsf{ledger}}_e=\hash(\operatorname{canon}(\Phi_e)).
\]
Because every non-genesis block contains the hash of its predecessor, each frontier transitively
commits to every block on its account chain from genesis through that frontier, including blocks
finalized in previous epochs.  Replay of those chains deterministically commits to balances, pending
and received sends, and representatives.  The preceding close-record hash makes the decided close
records a hash chain over the complete ledger history.

\paragraph{Disseminated and derived objects.}
Blocks, reports, and vote batches are immutable authenticated protocol objects.  A correct replica
advertises and transfers each unique valid such object only to peers that do not already possess it.  A
vote batch authenticates a canonical set of logical vote leaves; receivers expand valid leaves into the
same per-election state transitions that scalar votes would produce.  A certificate is not a wire
object: it is the local fact that the expanded validated vote store contains a quorum.  Report
proofs, conflicts, close-round death facts, close versions, outcomes, decisions, frontiers, release
state, and committees are likewise local derivations.  Inventory advertisements, retrieval requests,
and reconciliation deltas may be exchanged, but they carry no consensus authority.

\paragraph{Validated-close abstraction.}
A correct replica votes for a close hash only after it has locally constructed or reconstructed the
corresponding canonical set or frontier map and validated it using the underlying blocks, reports, and
votes.  Fresh admissibility is evaluated against the replica's current validated evidence and its own
signed-report and ordinary-finality locks before a new close vote.  Once a close-cut or close-record hash
obtains persistent same-hash certificate support in every applicable committee, its hash-checked
canonical preimage is historically admissible and never loses validity; later report equivocation or
newly learned evidence affects only fresh construction.  For a close record, the quorum-intersection
argument below guarantees that a persistently supported hash cannot omit an ordinary-finalized
segment: the correct signer weight locked by the slot certificate is greater than \(F_Q+P_Q\), leaving
insufficient weight for incompatible close-record support.  Authority depends only on locally derived
vote quorums and the hash-checked canonical preimage, not on certificate or transport-message arrival
order.

After the close-cut decision, the authoritative cut is the version reconstructed for the decided hash.
The cut is supporting close-protocol state: it freezes the unresolved elections that may participate in
close drain, but it is not the durable ledger-consensus value and is not an additional field of the
close-record hash.  Ordinary-finalized epoch-\(e\) segments may advance the close-record frontier even
when their elections are outside the cut.  After the close-record decision, the durable consensus
checkpoint is \(\closeHash[e]=h_F\), while \(\closeState[e]\) is reconstructed from the decided frontier
map, the decided cut, and the underlying validated blocks, reports, and vote batches.  Only this close
result, not local drain observations, feeds committee derivation and future slot-voting guards.  Each election id has its own timeout value
\(\timeoutVal_\id\).  In particular, \(\timeoutVal_{\Slot(a,s,e)}\) does not create an empty segment,
finalize any block, or skip the starting account slot \((a,s)\).  A later retry
\(\Slot(a,s,e')\) is permitted only after the certified close state releases the earlier election; a
segment finalized by an ordinary certificate or the close record is not retried.

For a candidate tip \(b=\hash(B_t)\) in \(q=\Slot(a,s,e)\), a replica stores every known block
\(B_r\) on the candidate suffix as
\begin{center}
\begin{adjustbox}{max width=\columnwidth}
\(\displaystyle
    \blockMap[q,\hash(B_r)]=B_r,\qquad s\le r\le t.
\)
\end{adjustbox}
\end{center}
The tip entry \(\blockMap[q,b]=B_t\) identifies the voted value; predecessor hashes identify the
remaining segment blocks.
Uninitialized map entries have the following defaults:
\begin{center}
\begin{adjustbox}{max width=\columnwidth}
\(\displaystyle
\begin{gathered}
    epochStart[e]=\bot,
    \quad T_{\mathrm{election}}[\id]=\bot,
    \quad outcome[\id]=\pending,\\
    \closeDecision[I]=\bot,
    \qquad \closeRound[I]=0,
    \quad I\in\{\CloseCut(e),\CloseRecord(e)\},\\
    notarized[P_i,Q,\id]=secondLook[P_i,Q,\id]=\{\},
    \quad firstVote[P_i,\id,Q]=\bot,
    \quad finalVoted[P_i,\id,Q]=\false,\\
    \votedFor[P_i,\id,b]=\false,
    \quad \blockMap[\cdot,\cdot]=\bot,
    \quad \committee[\cdot]=\bot,
    \quad \closeHash[e]=\bot,
    \quad \ledgerRoot[e]=\bot,
    \quad \closeState[e]=\bot,
    \quad \cutVersion[e]=\{\},
    \quad \frontierVersion[e]=\{\},
    \quad \carry[\id]=\bot,
    \quad \closeCut[e]=\bot.
\end{gathered}
\)
\end{adjustbox}
\end{center}
Account-slot finality is tracked by
\begin{center}
\begin{adjustbox}{max width=\columnwidth}
\(\displaystyle
    \finalizedSlot[a,s]\in\{\bot\}\cup
    \{b:\exists B,\ b=\hash(B),\ \slotOf(B)=(a,s)\},
\)
\end{adjustbox}
\end{center}
initially \(\finalizedSlot[\cdot,\cdot]=\bot\).  A block is finalized exactly when its account slot
maps to its hash in \(\finalizedSlot\); no separate block-finality map is maintained.  Because finality
is gap-free on each account chain, define
\[
\NextSlot(a)=1+\max\bigl(\{0\}\cup\{s:\finalizedSlot[a,s]\neq\bot\}\bigr).
\]
An enabled segment election for account \(a\) must start at \(\NextSlot(a)\).

Derived predicates and outcomes are
\begin{center}
\begin{adjustbox}{max width=\columnwidth}
\(\displaystyle
    \started[\id]\iff T_{\mathrm{election}}[\id]\neq\bot,
    \qquad
    \firstVoted(P_j,\id,Q)\iff firstVote[P_j,\id,Q]\neq\bot,
    \qquad
    \finalVoted(P_j,\id,Q)\iff finalVoted[P_j,\id,Q]=\true,
    \qquad
    \finished(\id)\iff outcome[\id]\neq\pending,
\)
\end{adjustbox}
\end{center}
\begin{center}
\begin{adjustbox}{max width=\columnwidth}
\(\displaystyle
    outcome[\id]\in
    \{\pending,\ \notarized(x),\ \converged(x),\ \timeoutOutcome,\ \fast(x),\ \final(x)\}.
\)
\end{adjustbox}
\end{center}
For a close instance \(I\), the variable
\[
    \closeDecision[I]\in
    \{\bot\}\cup\{(r,h)\}
\]
is instance-wide rather than round-local.  Here \(r\) is the deciding close round and \(h\) is the
decided hash.  The deciding certificate is derived locally from signed vote leaves, and the canonical
preimage is retained separately in \(\cutVersion\) or \(\frontierVersion\).

The round-local variable \(outcome[I_r]\) records only local progress in round \(r\).  It is not the
authoritative close decision and does not prevent a replica from processing a valid fast or final
certificate learned later as a decision.  Define
\[
    \closeFinished(I)\iff\closeDecision[I]\neq\bot.
\]
Local completion is represented by membership of every candidate block in \(\completeTree\), not by
an election outcome.  A complete segment is neither finished nor settled by completion alone.  A
non-timeout slot notarization certificate may finish a segment election as \(\notarized(b)\) for close
drain, but \(\notarized(b)\) does not finalize any block in the segment.  At close, each replica uses its currently known valid vote leaves and blocks to derive a fresh canonical
frontier map.  Global fast or final evidence is ordinary finality and may advance the epoch-\(e\) frontier
whether or not the election is in the close cut.  For elections in the cut, locally derived timeout or
conflicting notarizations propose omission and release, while exactly one currently known non-timeout
notarization, with no known final or release evidence, may propose \(b\) for inclusion on a candidate
account chain ending at the proposed frontier.  Candidate inclusion is not finality.  The complete
segment committed by tip \(b\) finalizes through ordinary fast/final evidence or when the close-record
hash committing to that frontier map is stored in \(\closeDecision[\CloseRecord(e)]\) with a valid fast
or final certificate.  A decided frontier that omits \(b\) releases \(b\); an ordinary-finalized tip
cannot be omitted by a valid close-record certificate because the report/finality lock intersection
prevents incompatible close-record support.

A \(\FirstCert(\Slot(a,s,e),b)\), \(\NotarCert(\Slot(a,s,e),b)\), and
\(\FinalCert(\Slot(a,s,e),b)\) are compatible certificates for the same segment tip \(b\).  A fast or final
certificate for \(b\) is incompatible with a notarization certificate for any \(b'\neq b\), including
the timeout value.  Conversely, two conflicting notarization certificates prove that no value can now
or later obtain an ordinary global fast or final certificate for that election; before close
certification they support release, while after a selecting close certificate they are losing
non-final evidence.  If a replica records an election outcome for the first time, its replica-local
\(outcome[\id]\) changes once from \(\pending\) to a non-pending value.  This variable is not an
agreed consensus value: correct replicas may temporarily record different terminal observations for
the same election, and no safety-critical retry or finality decision is derived from it without the
corresponding verifiable certificates and certified close state.  Later fast or final slot certificates
can still finalize their complete segment before the close record decides the obligation.  For
\(y\in\{\notarized(x),\fast(x),\final(x)\}\), \(\valueOf(y)=x\).  Election vote state is
logical committee-local state: \(firstVote[P_i,\id,Q]\), \(finalVoted[P_i,\id,Q]\),
\(notarized[P_i,Q,\id]\), and \(secondLook[P_i,Q,\id]\) are maintained independently for every
\(Q\in\Committees(\id)\).  Define the complete signed support of \(P_i\) in that committee instance by
\[
\func{Support}(P_i,Q,\id)=
notarized[P_i,Q,\id]\cup\bigl(\{firstVote[P_i,\id,Q]\}\setminus\{\bot\}\bigr).
\]
The first vote, including a timeout first vote, is therefore part of the support checked by every
final-vote rule.  Define the committees in which signer \(P_i\)'s vote has effect by
\[
\SignerCommittees(P_i,\id)=\{Q\in\Committees(\id):P_i\in Q\}.
\]
A logical vote carries an optional committee scope \(\chi\).  The value \(\chi=\bot\) means that the
scope field is omitted and the vote applies to every committee containing the signer.  An explicit
scope is a canonical nonempty proper subset of those committees.  Define
\begin{center}
\begin{adjustbox}{max width=\columnwidth}
\(\displaystyle
\EffectiveCommittees(P_i,\id,\chi)=
\begin{cases}
\SignerCommittees(P_i,\id), & \chi=\bot,\\
\chi, & \emptyset\neq\chi\subsetneq\SignerCommittees(P_i,\id).
\end{cases}
\)
\end{adjustbox}
\end{center}
The canonical sender convention omits the scope whenever the target set is all of
\(\SignerCommittees(P_i,\id)\); an explicit scope is used only for a proper subset.  A correct signer
emits an action into a nonempty target set \(S\) only when the action's committee-local guard holds in
every \(Q\in S\).  In particular, it may final-vote \(x\) into \(S\) only when
\[
\forall Q\in S:\quad
\func{Support}(P_i,Q,\id)\subseteq\{x\}\wedge\neg\finalVoted(P_i,\id,Q).
\]  Different first-vote values may be signed only for disjoint committee scopes; different
first-vote values with overlapping effective scopes constitute equivocation in every overlap.
Notarization plurality remains governed independently by the existing second-look and timeout rules in
each committee.  Upon sending or receiving a valid vote, a replica applies the corresponding state
transition separately and atomically to exactly the committees in its effective scope.


\paragraph{Vectorized signed vote transport.}
A logical vote leaf is
\[
\VoteLeaf=(\tau,\id,x,\chi),
\]
where \(\tau\in\{\mathsf{first},\mathsf{notar},\mathsf{final}\}\), \(x\) is the logical vote value,
and \(\chi\) is either omitted on the wire, represented here by \(\bot\), or is a canonical explicit
scope.  The election identifier contains its epoch and election type, from which every verifier derives
\(\Committees(\id)\), \(\SignerCommittees(P_i,\id)\), and then
\(\EffectiveCommittees(P_i,\id,\chi)\).  Because \(\chi\) is part of the canonical leaf, the batch
signature binds the scope.  For one signer and one governing close hash, let \(M\) be a canonical
duplicate-free set of leaves and let \(r_M=\hash(\operatorname{canon}(M))\).  The authenticated wire
object is
\[
\VoteBatch(P_i,\governingHash(e),r_M,M,\sigma_i),
\qquad
\sigma_i=\operatorname{Sign}_{P_i}
 \bigl(\hash(\mathsf{RAI/VoteBatch/v3}\parallel\governingHash(e)\parallel r_M)\bigr).
\]
A sender may transmit the complete batch or one leaf with a Merkle inclusion proof against \(r_M\).
A receiver validates the batch signature, canonical root, election identifier, governing context,
admissibility, committee derivation, scope canonicality, signer membership, and all ordinary
per-committee guards before expanding a leaf into one logical committee-local vote for every
\(Q\in\EffectiveCommittees(P_i,\id,\chi)\).  Duplicate leaves or multiple batches from one signer
contribute weight at most once to each \((\tau,\id,Q,x)\), and a batch cannot bypass committee-local
first-vote uniqueness, support restrictions, post-final locks, or certified release.  Logical votes for
independent account segments can therefore share one signature and transport header without sharing a
certificate or atomic commit domain.

Let \(G\) be the genesis epoch state.  Replay of \(G\) supplies genesis balances and
representatives.  For any certified ledger state through epoch \(k\), define delegated stake
\[
\stake_k(P_i)=
\sum_{\substack{a\in\mathsf{Accounts}\\\repOf_k[a]=P_i}}\bal_k[a].
\]
The exact-proportional committee snapshot is the weighted replica map
\[
Q_k=\{(P_i,\stake_k(P_i)):\stake_k(P_i)>0\},
\qquad
\weight_{Q_k}(P_i)=\stake_k(P_i).
\]
Thus every positive-stake representative is a committee member and its consensus weight is exactly the
total finalized balance delegated to it.  For any signer set \(S\subseteq Q\), write
\[
\SignerWeight_Q(S)=\sum_{P_i\in S}\weight_Q(P_i),
\qquad W_Q=\SignerWeight_Q(Q).
\]
Committee membership and weights are frozen for the lifetime of the snapshot.  Let
\begin{center}
\begin{adjustbox}{max width=\columnwidth}
\(\displaystyle
    \committee_{\mathsf{gen}}:=\deriveCommittee(\mathsf{genesis},G)=Q_{\mathsf{gen}}.
\)
\end{adjustbox}
\end{center}
Each closed epoch derives
\begin{center}
\begin{adjustbox}{max width=\columnwidth}
\(\displaystyle
    \committee[e]:=\deriveCommittee(e,\closeHash[e],\closeState[e])=Q_e.
\)
\end{adjustbox}
\end{center}
The close hash identifies the certified snapshot, while replay of \(\closeState[e]\) supplies the
balances and representatives from which the weight map is deterministically reconstructed.
The certified close state that participates in every election of epoch \(e\) is identified by
\begin{center}
\begin{adjustbox}{max width=\columnwidth}
\(\displaystyle
\governingHash(e)=
\begin{cases}
\hash(G), & e<2,\\
\closeHash[e-2], & e\ge 2.
\end{cases}
\)
\end{adjustbox}
\end{center}
For \(e\ge2\), a replica may validate or emit an epoch-\(e\) vote only after it has obtained and
validated \(\closeState[e-2]\) and its close certificate.  This governing context is implicit: it is
not a serialized vote field and is not separately signed.  A verifier derives
\(\governingHash(e)\) from its certified close history and evaluates the vote, close-version data, and
committee lookup under that context; if the required certified close state is unavailable or
inconsistent, the vote is invalid for that verifier.
Committee lookup is
\begin{center}
\begin{adjustbox}{max width=\columnwidth}
\(\displaystyle
\committeeAt(k)=
\begin{cases}
\committee_{\mathsf{gen}}, & k<0,\\
\committee[k], & k\ge 0 \text{ and } \committee[k]\neq\bot,\\
\bot, & k\ge 0 \text{ and } \committee[k]=\bot.
\end{cases}
\)
\end{adjustbox}
\end{center}
The deterministic reporting committee for epoch \(e\) is
\begin{center}
\begin{adjustbox}{max width=\columnwidth}
\(\displaystyle
\ReportCommittee(e)=\committeeAt(e-2).
\)
\end{adjustbox}
\end{center}
It is used for the close-report barrier, report-derived visibility, and both close-election phases.
Slot finality continues to use every slot-election committee below.

Election committees are
\begin{center}
\begin{adjustbox}{max width=\columnwidth}
\(\displaystyle
    \Committees(\Slot(a,s,e))=\{\committeeAt(e-3),\committeeAt(e-2)\},
    \qquad
    \Committees(\CloseCut(e,r))=\{\committeeAt(e-2)\},
    \qquad
    \Committees(\CloseRecord(e,r))=\{\committeeAt(e-2)\},
\)
\end{adjustbox}
\end{center}
with duplicate committees collapsed.  Each \(Q\in\Committees(\id)\) runs the election core
independently.  A vote for \(\id\) is valid only when
\(\SignerCommittees(P_i,\id)\neq\varnothing\), every committee lookup used by
\(\Committees(\id)\) is non-\(\bot\), its effective scope is well formed and nonempty, and the signer
satisfies the applicable vote guard in every committee in
\(\EffectiveCommittees(P_i,\id,\chi)\).  An omitted scope is interpreted as all committees in
\(\SignerCommittees(P_i,\id)\); an explicit scope must be a canonical proper subset.  Every verifier
derives the election committees from \(\id\), tallies the signer once in each committee in the effective
scope, and does not tally it elsewhere.  The vote occupies the corresponding vote slot and any
post-final lock independently in each targeted committee; duplicate copies from one signer are
deduplicated per committee.  Every logical vote leaf commits to \(\id\), its value, and its optional
scope, while its governing close state remains implicit validation context derived from the vote's epoch
and certified close history.  This lagged committee lookup permits epoch \(e+1\) to open as soon as
epoch \(e\) enters closing: its governing close state
is epoch \(e-1\), which is already closed by the precondition of \(\StartClosing(e)\), while
\(\committee[e]\) is derived only when epoch \(e\) later finishes closing.

\paragraph{Performance profile.}
If a segment contains \(k\) blocks, one election round amortizes committee voting over those \(k\)
blocks.  If one signed batch contains \(m\) logical vote leaves, signature generation and fixed vote
headers are amortized over those \(m\) elections.  Validation and ledger application remain linear in
the number of blocks and operations, while independent account segments can be validated and certified
in parallel.  Cross-account dependency finality is deliberately excluded from this fast path: a receive
can consume only a send finalized before the segment base.

\section{Voting Guards and Release}

A replica may cast first votes, timeout first votes, extra notarization votes, and timeout
notarization votes only while the election is enabled:
\begin{center}
\begin{adjustbox}{max width=\columnwidth}
\(\displaystyle
\begin{aligned}
\enabled(\Slot(a,s,e))\iff{}& s=\NextSlot(a)\wedge\finalizedSlot[a,s]=\bot\\
&\wedge\bigl(state[e]=\mathsf{open}
\vee (state[e]=\mathsf{closing}\wedge\closeCut[e]\neq\bot\wedge\Slot(a,s,e)\in\closeCut[e])\bigr),\\
\enabled(\CloseCut(e,r))\iff{}& state[e]=\mathsf{closing},\\
\enabled(\CloseRecord(e,r))\iff{}& state[e]=\mathsf{closing}.
\end{aligned}
\)
\end{adjustbox}
\end{center}
Thus an account has at most one enabled unresolved segment obligation: the election beginning at its
next unfinalized slot.  A segment election is normally vote-enabled only while its epoch is open, but
after a certified close cut is installed it remains vote-enabled during close drain exactly when it is
included in that cut.

For a correct replica \(P_j\), casting any non-timeout first, notarization, or final vote for segment
tip \(b\) in \(\Slot(a,s,e)\) records an unreleased vote:
\begin{center}
\begin{adjustbox}{max width=\columnwidth}
\(\displaystyle
    \votedFor[P_j,\Slot(a,s,e),b]=\true.
\)
\end{adjustbox}
\end{center}
For \(q=\Slot(a,s,e)\), define replica-local ordinary finality and positive non-finalizability evidence by
\[
\begin{aligned}
\ordinaryFinalized_j(q,b)\iff{}&
 b\neq\timeoutVal_q\wedge \Complete(q,b)\wedge
 \bigl(\FirstCert(q,b)\in\mathsf{pool}_j\vee\FinalCert(q,b)\in\mathsf{pool}_j\bigr)\\
&\wedge\SegmentBlocks(q,b)=(B_s,\ldots,B_t)
 \wedge\forall r\in[s,t]:\finalizedSlot_j[a,r]=\hash(B_r),\\
\cannotFinalize_j(q)\iff{}&
 \text{the validated vote leaves known to }P_j\text{ derive timeout evidence or}\\
&\text{conflicting terminal-certificate evidence for }q.
\end{aligned}
\]
After \(P_j\) signs its epoch report \(R_j[e]\), define the fresh report lock
\[
\reportLocked_j(q)\iff q\in R_j[e]\wedge\neg\cannotFinalize_j(q).
\]
Learning ordinary finality does not erase this signed-report history; it replaces the unresolved
inclusion obligation by an exact lock on the finalized tip during fresh close-record voting.

An earlier vote is released only when the certified close record omits that candidate from the decided
ledger:
\begin{center}
\begin{adjustbox}{max width=\columnwidth}
\(\displaystyle
\begin{aligned}
\released_{P_j}(q,b)
\iff{}& q=\Slot(a,s,e)\wedge state[e]=\mathsf{closed}\wedge\closeCut[e]\neq\bot\\
&{}\wedge\neg\func{LedgerIncludes}_e(q,b).
\end{aligned}
\)
\end{adjustbox}
\end{center}
Define \(\func{LedgerIncludes}_e(q,b)\) to hold when \(q=\Slot(a,s,e)\),
\(\blockMap[q,b]=B_t\), and the complete segment \(\SegmentBlocks(q,b)=(B_s,\ldots,B_t)\)
occurs as the account-chain suffix committed by \(R^{\mathsf{ledger}}_e\).  If the certified ledger
contains a different segment from the same starting slot, every conflicting tip is released.  If no
candidate of \(q\) occurs on the certified ledger, the election is released, whether or not \(q\) was
in the close cut.  Cut omission alone is not release before the close record certifies.  Raw timeout or
conflict evidence may finish the election for drain but does not by itself release future conflicting
votes while the close record is unsettled.  A local belief that an election is stale, idle, or disabled
is not release.

The expanded certified close state maps an election \(q\) to
\(\FinalizedState(b,\pi_{\mathsf{ledger}})\) when \(\func{LedgerIncludes}_e(q,b)\), and otherwise to
\(\ReleasedState(\pi)\), with close-exclusion, timeout, or conflict evidence derived from the validated
report and vote store.  The state is cumulative with respect to finality: every previously finalized
block must occur on the certified account chains, and no later close state may replace it with a
conflicting value.  For account \(a\), define the governing certified frontier
\begin{center}
\begin{adjustbox}{max width=\columnwidth}
\(\displaystyle
\requiredFrontier(a,e)=
\begin{cases}
\frontier[e-2,a], & e\ge2\text{ and }\closeState[e-2]\text{ is certified},\\
\bot, & \text{otherwise.}
\end{cases}
\)
\end{adjustbox}
\end{center}
Write \(\DescendsFrom(B,b_c)\) when \(b_c\) occurs on the strict parent chain of \(B\).  For
\(q=\Slot(a,s,e)\) and candidate tip \(B_t\), define the certified-base lock
\begin{center}
\begin{adjustbox}{max width=\columnwidth}
\(\displaystyle
\StablePrefix(q,B_t)\iff
\bigl(s=1\wedge\parent(B_s)=\mathsf{genesis}_a\bigr)
\vee
\bigl(s>1\wedge\parent(B_s)=\finalizedSlot[a,s-1]\bigr)
\)
\end{adjustbox}
\end{center}
where \((B_s,\ldots,B_t)=\SegmentBlocks(q,\hash(B_t))\).  Then
\begin{center}
\begin{adjustbox}{max width=\columnwidth}
\(\displaystyle
\begin{aligned}
\safeFromClose(\Slot(a,s,e),b)\iff{}& \exists B_t:\ b=\hash(B_t)\wedge\blockMap[\Slot(a,s,e),b]=B_t\\
&\wedge\StablePrefix(\Slot(a,s,e),B_t)\wedge\SegmentValid(\Slot(a,s,e),b)\\
&\wedge\bigl(e<2\ \vee\ (\closeState[e-2]\text{ is certified}
\wedge\DescendsFrom(B_t,\requiredFrontier(a,e)))\bigr).
\end{aligned}
\)
\end{adjustbox}
\end{center}
The certified-base lock replaces the per-block finalized-ancestor bottleneck.  Blocks inside one
candidate segment may be unresolved because the certificate decides them together, but every ancestor
before the starting slot is already finalized and the entire segment is validated before voting.
Because \(s=\NextSlot(a)\), no later slot election for the same account is enabled while the segment is
unresolved.  During close drain, every non-final selected block beyond \(\frontier[e-1,a]\) must belong to the
selected segment of an election in \(\closeCut[e]\); an ordinary-finalized epoch-\(e\) segment may also
advance the close-record frontier outside the cut.  Every receive must resolve to a send finalized
before that segment's base snapshot.  A released segment can never satisfy the base lock for a
descendant.  A certified close record freezes every epoch-\(e\) election as finalized or released;
future votes for the account may extend only the certified finalized frontier.

Earlier safety-relevant vote history for account slot \((a,s)\) is
\begin{center}
\begin{adjustbox}{max width=\columnwidth}
\(\displaystyle
\begin{aligned}
\activeVote_{P_j}(a,s,e_0,b)
\iff{}& \votedFor[P_j,\Slot(a,s,e_0),b]=\true.
\end{aligned}
\)
\end{adjustbox}
\end{center}
A later non-timeout slot vote is safe exactly when
\begin{center}
\begin{adjustbox}{max width=\columnwidth}
\(\displaystyle
\begin{aligned}
\safeToVote_{P_j}(\Slot(a,s,e'),b')
\iff{}& s=\NextSlot(a)\wedge\finalizedSlot[a,s]=\bot \\
&\wedge\safeFromClose(\Slot(a,s,e'),b')\\
&\wedge\forall e_0<e',\ \forall b\neq b':
\activeVote_{P_j}(a,s,e_0,b)\Rightarrow
\released_{P_j}(\Slot(a,s,e_0),b).
\end{aligned}
\)
\end{adjustbox}
\end{center}
Thus a correct replica may vote for different segment tips in disjoint committee instances of the same
election as allowed by the election core, while retaining at most one first vote and one post-final lock
per committee.  It may vote in a later retry of the same starting slot only after every earlier
obligation has been certifiably released.  Once a segment finalizes, \(\NextSlot(a)\) advances past its
tip, so internal slots are never opened as competing elections.  Later segments must descend from the
governing certified frontier and begin at the current next slot.  Timeout and close votes do not create
slot unreleased votes.

\section{Certificates and Vote Counting}

RAI instantiates, for each fixed election identifier and weighted committee, the vote-counting core of
Kudzu~\cite{shoup2025kudzu}.  Kudzu supplies the single-committee counting pattern behind its
fast-path safety lemmas and its two-case termination argument.  RAI replaces voter cardinality with the
frozen consensus weight of the applicable stake snapshot.  It does not import Kudzu's block proposal,
reconstruction, or slot-advance semantics.  Optional-scope deterministic vote expansion,
election-scoped admissibility, multiple-committee composition, close-round carry, close-version
reconciliation, and cross-epoch reasoning are RAI-specific and are stated and proved below.

\paragraph{Committee-local election kernel.}
Fix an election identifier \(\id\) and a weighted committee \(Q\in\Committees(\id)\).  Every vote
signature binds the vote kind, \(\id\), value, and optional scope; the signer identity, \(\id\), and
scope deterministically derive every committee-local instance in which the vote is counted.  A correct
member sends at most one first vote for each \((\id,Q)\), and a first-vote leaf counts as both a first
and notarization vote in every committee in its effective scope.  A correct member may issue a
second-look notarization vote for a non-timeout value \(x\) into a target set \(S\) only after validating
\(\admissible(\id,x)\) and observing first-vote signer weight strictly greater than
\(F_{Q'}+P_{Q'}\) for \(x\) in every \(Q'\in S\).  It may issue a timeout notarization vote into
\(S\) only when \(\timeoutReady(Q',\id)\) holds for every \(Q'\in S\).  Before final-voting \(x\)
into \(S\), its complete committee-local support is a subset of \(\{x\}\) in every \(Q'\in S\).
After final-voting in a committee, it emits no later first, notarization, timeout, or final vote in that
committee instance, but it may continue with independently valid actions in disjoint committee scopes.
Valid votes and certificates are persistent.  These rules are the complete committee-local assumptions
used when adapting Kudzu's Lemmas~5.4--5.7 and Lemma~5.10; no other property of Kudzu is assumed.
Notation such as \(\NotarVote(P_i,\id,Q,x)\) denotes the derived \(Q\)-local expansion of a valid
signed leaf whose effective scope contains \(Q\), not a committee field added to the local vote.

Each committee \(Q\in\Committees(\id)\) has total weight \(W_Q\), maximum Byzantine weight
\(F_Q\), and quorum-slack weight \(P_Q\).  Assume
\[
P_Q\ge F_Q,
\qquad
W_Q>3F_Q+2P_Q.
\]
A local non-timeout certificate \(\Cert_Q(k,\id,x)\), where
\(k\in\{\notarKind,\fastKind,\finalKind\}\) and \(x\neq\timeoutVal_\id\), has threshold
\begin{center}
\begin{adjustbox}{max width=\columnwidth}
\(\displaystyle
\begin{array}{c|c|c}
\text{kind} & \text{signer-weight threshold in }Q & \text{value restriction}\\
\hline
\notarKind & \ge W_Q-F_Q-P_Q \text{ notarization weight} & \text{none}\\
\fastKind & \ge W_Q-P_Q \text{ first-vote weight} & x\neq\timeoutVal_\id\\
\finalKind & \ge W_Q-F_Q-P_Q \text{ final-vote weight} & x\neq\timeoutVal_\id.
\end{array}
\)
\end{adjustbox}
\end{center}
A certificate is the local fact that individually signed votes of sufficient distinct-signer weight are
present; it is not a separate BLS, threshold-signature, or extra vote-bearing object.  All votes counted
for a local certificate are valid for the same \((\id,Q,x)\), and each signer contributes its frozen
committee weight at most once.  A first vote for \((\id,Q,x)\) is also counted as a notarization vote
for that tuple.  Valid votes and certificates do not expire.

A local timeout certificate \(\TimeoutCert(\id,Q)\) exists when valid timeout notarization votes have
signer weight at least \(W_Q-F_Q-P_Q\).  The local terminal-result domain is
\[
L_Q(\id)\in\{\timeoutOutcome,\notarized(x),\fast(x),\final(x)\}.
\]
The RAI election result is the deterministic merge of the local results.  For the same non-timeout
value \(x\),
\begin{center}
\begin{adjustbox}{max width=\columnwidth}
\(\displaystyle
\begin{aligned}
\func{merge}(\notarized(x),\notarized(x))&=\notarized(x),\\
\func{merge}(\notarized(x),\fast(x))&=\func{merge}(\notarized(x),\final(x))=\notarized(x),\\
\func{merge}(\fast(x),\fast(x))&=\fast(x),\\
\func{merge}(\fast(x),\final(x))&=\func{merge}(\final(x),\final(x))=\final(x),\\
\func{merge}(\timeoutOutcome,y)&=\timeoutOutcome.
\end{aligned}
\)
\end{adjustbox}
\end{center}
The equations are symmetric.  Two non-timeout local results for different values, including two local
final results for different values, merge to \(\timeoutOutcome\).  Their certificate-level evidence is a
conflict proof that the round cannot produce a compatible global fast or final result.  With one
committee, the merged result is that committee's local result.  The global aliases
\(\NotarCert(\id,x)\), \(\FirstCert(\id,x)\), and \(\FinalCert(\id,x)\) mean that the merged result
is respectively \(\notarized(x)\), a joint fast result \(\fast(x)\) from every election committee, or
\(\final(x)\); \(\NotarCert(\id,\timeoutVal_\id)\) means that the merged result is
\(\timeoutOutcome\).

A merged result finishes an election round according to the slot/close terminal rules below.
Fast and final results are valid only for non-timeout values.  A non-timeout notarization result in the
unique close committee records round-local progress as \(\converged(x)\) and creates a live
carry unless positive death evidence is known.  A valid fast or final certificate in any close round
decides the logical close instance through \(\closeDecision\), regardless of the replica's current
round or the round-local value of \(outcome\).  Correct replicas continue to use the committee-local
election kernel above in close rounds.

\begin{center}
\begin{adjustbox}{max width=\columnwidth}
\(\displaystyle
\Terminal(t,\id,x)\iff
\begin{cases}
    t\in\{\notarKind,\fastKind,\finalKind\}, & \id=\Slot(a,s,e),\\
    t=\notarKind\vee(t\in\{\fastKind,\finalKind\}\wedge x\neq\timeoutVal_\id),
        & \id\in\{\CloseCut(e,r),\CloseRecord(e,r)\},
\end{cases}
\)
\end{adjustbox}
\end{center}
where \(t,x\) denote the kind and value represented by the merged result, and
\begin{center}
\begin{adjustbox}{max width=\columnwidth}
\(\displaystyle
\certOutcome(t,x)=
\begin{cases}
\timeoutOutcome, & t=\notarKind\wedge x=\timeoutVal_\id,\\
\notarized(x), & t=\notarKind\wedge x\neq\timeoutVal_\id\wedge\id=\Slot(a,s,e),\\
\converged(x), & t=\notarKind\wedge x\neq\timeoutVal_\id\wedge\id\in\{\CloseCut(e,r),\CloseRecord(e,r)\},\\
\fast(x), & t=\fastKind,\\
\final(x), & t=\finalKind.
\end{cases}
\)
\end{adjustbox}
\end{center}

For \(Q\in\Committees(\id)\), define
\begin{center}
\begin{adjustbox}{max width=\columnwidth}
\(\displaystyle
\VotesIn(Q,\id,x)=\{P_i\in Q: firstVote[P_i,\id,Q]=x\},
\qquad
\NotarVotesIn(Q,\id,x)=
\{P_i\in Q: firstVote[P_i,\id,Q]=x\vee\NotarVote(P_i,\id,Q,x)\in\mathsf{pool}\}.
\)
\end{adjustbox}
\end{center}
\begin{center}
\begin{adjustbox}{max width=\columnwidth}
\(\displaystyle
\allVotes(Q,\id)=
\SignerWeight_Q(\{P_i\in Q:firstVote[P_i,\id,Q]\neq\bot\}),
\)
\end{adjustbox}
\end{center}
\begin{center}
\begin{adjustbox}{max width=\columnwidth}
\(\displaystyle
\maxVotes(Q,\id)=\max_{x\neq\timeoutVal_\id}\SignerWeight_Q(\VotesIn(Q,\id,x)),
\)
\end{adjustbox}
\end{center}
where the maximum is \(0\) if there is no non-timeout value, and
\begin{center}
\begin{adjustbox}{max width=\columnwidth}
\(\displaystyle
\manyVotes(Q,\id)=
\{x\neq\timeoutVal_\id:\ \SignerWeight_Q(\VotesIn(Q,\id,x))>F_Q+P_Q\}.
\)
\end{adjustbox}
\end{center}
Timeout notarization for committee \(Q\) is enabled by
\begin{center}
\begin{adjustbox}{max width=\columnwidth}
\(\displaystyle
\timeoutReady(Q,\id)\iff
\allVotes(Q,\id)-\maxVotes(Q,\id)>F_Q+P_Q.
\)
\end{adjustbox}
\end{center}

The monotone slack rule above is the signed timeout trigger, following Kudzu's main loop
\cite{shoup2025kudzu}.  For diagnostics, let \(\Values(\id)\) be the finite election-local set of
non-timeout values that currently appear in valid close-version data or valid votes for \(\id\), and
define
\begin{center}
\begin{adjustbox}{max width=\columnwidth}
\(\displaystyle
\NotarizedBy(P_i,\id,Q)=
\{x:firstVote[P_i,\id,Q]=x\vee\NotarVote(P_i,\id,Q,x)\in\mathsf{pool}\}.
\)
\end{adjustbox}
\end{center}
\begin{center}
\begin{adjustbox}{max width=\columnwidth}
\(\displaystyle
\begin{aligned}
\FastDead(Q,\id,y)\iff{}&
\SignerWeight_Q(\{P_i\in Q:firstVote[P_i,\id,Q]\notin\{\bot,y\}\})>F_Q+P_Q,\\
\TornFor(Q,\id,y)={}&
\{P_i\in Q:\exists x\in\NotarizedBy(P_i,\id,Q),\ x\neq y\}
\setminus\{P_i\in Q:\FinalVote(P_i,\id,Q,y)\in\mathsf{pool}\},\\
\FinalDead(Q,\id,y)\iff{}&
\SignerWeight_Q(\TornFor(Q,\id,y))>2F_Q+P_Q,\\
\AllValuesDead(Q,\id)\iff{}&
\forall y\in\Values(\id):\FastDead(Q,\id,y)\wedge\FinalDead(Q,\id,y).
\end{aligned}
\)
\end{adjustbox}
\end{center}
\begin{center}
\begin{adjustbox}{max width=\columnwidth}
\(\displaystyle
\timeoutAllowed(Q,\id)\iff \timeoutReady(Q,\id).
\)
\end{adjustbox}
\end{center}
The set \(\Values(\id)\) is learned incrementally, so \(\AllValuesDead(Q,\id)\) is only a local
diagnostic over currently known values.  It does not authorize a signed timeout vote: an unknown value
could later acquire first-vote weight, and a timeout vote based on an incomplete candidate set would not
be covered by the weighted incompatibility argument.  A correct replica may use \(\AllValuesDead\) to
avoid futile local reconstruction after certificate-level death evidence is already present, but a
timeout notarization vote is issued only under the monotone \(\timeoutReady\) predicate.  Such a vote is
an ordinary signed notarization vote for \(\timeoutVal_\id\); receivers validate and count it like any
other notarization vote.

\paragraph{Lemma (committee-local certificate incompatibility).}
Fix \((\id,Q)\).  Assume Byzantine signers have total weight at most \(F_Q\) and
\(W_Q>3F_Q+2P_Q\).  If a valid local fast or final certificate exists for \(x\), then no valid local
notarization certificate exists for any \(y\neq x\), including \(\timeoutVal_\id\).

For the fast case, a fast certificate for \(x\) contains first-vote signer weight at least
\(W_Q-P_Q\).  Even allowing all Byzantine weight to equivocate, at most \(F_Q+P_Q\) distinct-signer
weight can be recorded outside \(x\).  Hence no different value can enter \(\manyVotes(Q,\id)\), whose
threshold is strictly greater than \(F_Q+P_Q\), and
\(\allVotes(Q,\id)-\maxVotes(Q,\id)\le F_Q+P_Q\), so \(\timeoutReady(Q,\id)\) is false.  Correct
members therefore send no second-look or timeout notarization vote for a value different from \(x\).
The remaining Byzantine weight is less than the notarization threshold \(W_Q-F_Q-P_Q\).

For the final case, a final certificate for \(x\) and a notarization certificate for \(y\neq x\) would
have signer weights at least \(W_Q-F_Q-P_Q\).  Their intersection has weight at least
\[
2(W_Q-F_Q-P_Q)-W_Q=W_Q-2F_Q-2P_Q>F_Q.
\]
It therefore contains correct signer weight, contradicting the support condition and post-final lock of
the committee-local kernel.  This is the stake-weighted RAI instantiation of Kudzu's fast-finalization
and incompatibility arguments~\cite[Lemmas~5.4--5.7]{shoup2025kudzu}.

Since a global fast or final certificate requires the same value in every election committee, any two
valid committee-local non-timeout terminal certificates for different values in applicable committees
form a conflict proof
\[
\ConflictProof(\id,x,y),\qquad x\neq y,
\]
where a terminal certificate may be notarization, fast, or final evidence.  Such evidence excludes every
possible present or future compatible global fast or final certificate for \(\id\).  A merged timeout
result records either an explicit timeout certificate or such incompatible terminal evidence; the
underlying signed vote leaves remain in the validated object store as the release witness.

\paragraph{Lemma (committee composition).}
Suppose every \(Q\in\Committees(\id)\) has derived a local terminal result.  Every correct verifier
derives the same effective scope and committee-local expansion of each signed vote, while thresholds,
signer weights, certificates, and locks are evaluated per \((\id,Q)\).  The targeted-committee emission
guard makes every correct vote valid in each committee where it is counted, and committee-local
uniqueness prevents a correct signer from supporting different values inside the same committee even
when it signs different values for disjoint scopes.  The deterministic merge therefore produces a
global terminal result immediately.  It finalizes a non-timeout value only when every committee supplies
a compatible fast or final result for that same value.  Timeout in one committee, or different
non-timeout values in two committees, cannot produce finality.  By the preceding lemma, two different
globally finalized values are impossible.

\subsection{Close-round carry and death evidence}

Fix a close instance
\[
I\in\{\CloseCut(e),\CloseRecord(e)\},
\]
and write \(I_r\) for its election identifier in close round \(r\).

A close-round death proof is positive certificate evidence that round \(r\) cannot now or later
produce a deciding fast or final certificate:
\[
\CloseDead_{P_j}(I_r)\iff
\begin{cases}
\NotarCert(I_r,\timeoutVal_{I_r})\in\mathsf{pool}_{P_j},\\
\text{or }P_j\text{ has }\CloseConflictProof(I_r,h,h')\text{ for }h\neq h'.
\end{cases}
\]

A conflict proof consists of valid certificate-level committee-local terminal evidence, including
notarization, fast, or final evidence, whose coexistence is incompatible with every possible deciding
fast or final certificate for \(I_r\).  The underlying signed vote leaves are retained and disseminated through ordinary vote-batch gossip.  A local clock timeout, changed preference, or newly learned
uncertified evidence is not a death proof.

For a non-timeout hash \(h\), define
\[
\LiveCarry_{P_j}(I_r,h)\iff
\NotarCert(I_r,h)\in\mathsf{pool}_{P_j}
\wedge
\neg\CloseDead_{P_j}(I_r).
\]

The predicate is local.  Correct replicas may temporarily derive different live carries if a
the authenticated vote leaves completing a conflicting quorum have not yet been completely gossiped.  Such disagreement is safe because the
hidden conflict makes the source round incapable of deciding.

\section{Blocks, Admissibility, and Finality}

The complete block tree \(\completeTree\) contains block data, with one account tree per account.  A
complete account block is represented as \(B\), and \(\slotOf(B)=(a,s)\).  Define
\begin{center}
\begin{adjustbox}{max width=\columnwidth}
\(\displaystyle
\CompleteBlock(b)\iff
b\in\{\mathsf{genesis}_a:a\in\mathsf{Accounts}\}\vee\exists B\in\completeTree:b=\hash(B).
\)
\end{adjustbox}
\end{center}
\begin{center}
\begin{adjustbox}{max width=\columnwidth}
\(\displaystyle
\begin{aligned}
\ParentComplete(B,a,s;q)
\iff{}& (\parent(B)=\mathsf{genesis}_{a}\wedge s=1)\\
&\vee\exists B_p\in\bigl(\completeTree\cup\{\blockMap[q,h]:h\in\dom(\blockMap[q,\cdot])\}\bigr):
\parent(B)=\hash(B_p)\wedge\slotOf(B_p)=(a,s-1).
\end{aligned}
\)
\end{adjustbox}
\end{center}
Here \(\slotOf(\mathsf{genesis}_a)=(a,0)\).  Define
\[
\func{AncestorsConsistent}(B)\iff
\forall B_a\text{ on the parent chain of }B:\
\finalizedSlot[\slotOf(B_a)]\in\{\bot,\hash(B_a)\}.
\]

For \(q=\Slot(a,s,e)\) and tip \(b=\hash(B_t)\), define
\(\SegmentBlocks(q,b)=(B_s,\ldots,B_t)\) when these are the unique blocks on the parent chain of
\(B_t\) with account-local sequence numbers \(s,\ldots,t\), all are present in \(\blockMap[q,\cdot]\)
or \(\completeTree\), and the predecessor of \(B_s\) is the certified finalized block at slot
\(s-1\), or genesis when \(s=1\).

A receive reference \(u=(\hash(B_u),i)\) is eligible for a candidate segment only when \(B_u\) was
finalized before the segment's base snapshot, index \(i\) selects a send \((d,x)\), and the base
finalized ledger contains \(\pendingSend[u]=(d,x)\).  Sends created by blocks inside the same segment
are added to the temporary result but are not eligible to fund any receive in that segment.  This
preserves finalized send-before-receive dependencies, prevents cyclic or same-certificate cross-account
funding, and gives the segment one objective base ledger state.

Let \(L_0\) be that finalized base state and let \(L_r\) be temporary state after replaying through
\(B_r\).  For each block, debit and credit are
\[
\operatorname{credit}(B_r,L_{r-1},L_0)=
\sum_{u\in\receivesOf(B_r)}\pendingSend_{L_{r-1}}[u].\mathsf{amount},
\qquad
\operatorname{debit}(B_r)=\sum_{(d,x)\in\sendsOf(B_r)}x,
\]
with every receive additionally required to name an entry already present in
\(\pendingSend_{L_0}\).  The block-validity predicate in replay context is
\begin{center}
\begin{adjustbox}{max width=\columnwidth}
\(\displaystyle
\begin{aligned}
\ValidBlock(B_r;L_{r-1},L_0)\iff{}& \slotOf(B_r)=(a,r)\wedge\ParentComplete(B_r,a,r;q)\\
&\wedge\operatorname{Verify}_{\ownerKey[a]}
 \bigl(\signatureOf(B_r),
 \hash(\mathsf{RAI/AccountBlock/v1}\parallel\operatorname{canon}(\operatorname{body}(B_r)))\bigr)\\
&\wedge\text{every send amount is a positive representable integer}\\
&\wedge\text{all receive references in }B_r\text{ are distinct}\\
&\wedge\forall u\in\receivesOf(B_r):
 u\in\dom(\pendingSend_{L_0})\wedge u\notin\receivedSend_{L_{r-1}}\\
&\qquad\wedge\pendingSend_{L_{r-1}}[u]=(a,x)\text{ for some }x>0\\
&\wedge\bal_{L_{r-1}}[a]+\operatorname{credit}(B_r,L_{r-1},L_0)-\operatorname{debit}(B_r)\ge0\\
&\wedge\representativeOf(B_r)\in\mathsf{RegisteredReplicas}.
\end{aligned}
\)
\end{adjustbox}
\end{center}
Applying a valid block updates the temporary balance and representative, removes each consumed receive
from \(\pendingSend\), records it in \(\receivedSend\), and creates each new pending send.  Replay
also verifies the supply invariant.  Account-parent order is deterministic, while independent account
segments may be replayed in parallel because their receives depend only on previously finalized sends.

Define
\begin{center}
\begin{adjustbox}{max width=\columnwidth}
\(\displaystyle
\begin{aligned}
\SegmentValid(q,b)\iff{}&q=\Slot(a,s,e)\wedge b=\hash(B_t)\wedge
\SegmentBlocks(q,b)=(B_s,\ldots,B_t)\\
&\wedge t\ge s\wedge\func{AncestorsConsistent}(B_t)\wedge\StablePrefix(q,B_t)\\
&\wedge\forall r\in[s,t]:\finalizedSlot[a,r]\in\{\bot,\hash(B_r)\}\\
&\wedge\text{replay of }B_s,\ldots,B_t\text{ from }L_0\text{ satisfies }\ValidBlock
\text{ at every step and the supply invariant}.
\end{aligned}
\)
\end{adjustbox}
\end{center}
The complete block tree contains at most one entry for each block hash.  A slot value extends the tree
when
\begin{center}
\begin{adjustbox}{max width=\columnwidth}
\(\displaystyle
\ExtendsTree(\Slot(a,s,e),b)\iff
b\neq\timeoutVal_{\Slot(a,s,e)}\wedge\SegmentValid(\Slot(a,s,e),b).
\)
\end{adjustbox}
\end{center}
For a slot election \(q=\Slot(a,s,e)\), define a start witness by
\begin{center}
\begin{adjustbox}{max width=\columnwidth}
\(\displaystyle
\begin{aligned}
\func{StartWitness}_e(q)\iff{}&
\bigl(\exists Q\in\Committees(q):\mathsf{pool}\text{ contains a valid first vote for }q\text{ in }Q\bigr)\\
&\vee\Bigl(\SignerWeight_{\ReportCommittee(e)}
\bigl(\{P_i\in\ReportCommittee(e):q\in report[P_i,e]\}\bigr)>F_{\ReportCommittee(e)}\Bigr).
\end{aligned}
\)
\end{adjustbox}
\end{center}
The witness is a local predicate over already gossiped reports and votes; it is not a separate object.

\paragraph{Stable admissibility and fresh preference.}
Admissibility is distinct from the value that a correct replica currently prefers to first-vote.  For
correct replica \(P_j\), let \(E_j\) be its validated blocks, reports, and vote batches.  Canonical fresh
construction must satisfy
\[
E_i=E_j
\Longrightarrow
\FreshCloseHash_{P_i}(I)=\FreshCloseHash_{P_j}(I).
\]
Serialization order, witness choice, reconciliation path, and other transport choices do not affect a
close hash.

Slot admissibility remains
\begin{center}
\begin{adjustbox}{max width=\columnwidth}
\(\displaystyle
\admissible(\Slot(a,s,e),b)\iff
b\neq\timeoutVal_{\Slot(a,s,e)}\wedge\ExtendsTree(\Slot(a,s,e),b)
\wedge\safeFromClose(\Slot(a,s,e),b).
\)
\end{adjustbox}
\end{center}

A reconstructed close-cut version \(C\) is freshly admissible for \(h_C\) at correct replica \(P_j\)
when
\begin{center}
\begin{adjustbox}{max width=\columnwidth}
\(\displaystyle
\begin{aligned}
\admissibleFresh_{P_j}(\CloseCut(e),h_C,C)\iff{}&
 h_C=\hash(\mathsf{RAI/CloseCut}\parallel e\parallel\operatorname{canon}(C))\\
&\wedge\exists\rho:\ReportProof(e,\rho)\\
&\wedge\{\Slot(a,s,e):(a,s)\in\ReportVisible(e,\rho)\}\subseteq C\\
&\wedge\{q:q=\Slot(a,s,e)\wedge\reportLocked_j(q)\}\subseteq C\\
&\wedge\forall q\in C:\func{StartWitness}_e(q).
\end{aligned}
\)
\end{adjustbox}
\end{center}
Define persistent certificate support for a cut hash by
\begin{center}
\begin{adjustbox}{max width=\columnwidth}
\(\displaystyle
\begin{aligned}
\func{CutSupported}(e,h_C)\iff{}&\exists r:\ \forall Q\in\Committees(\CloseCut(e,r)),\\
&\exists k_Q\in\{\notarKind,\fastKind,\finalKind\}:\
\Cert_Q(k_Q,\CloseCut(e,r),h_C).
\end{aligned}
\)
\end{adjustbox}
\end{center}
This predicate is monotone because the underlying committee-local certificates and authenticated vote leaves
do not expire, even if later incompatible evidence changes the currently merged round result.
Historical cut admissibility and the combined predicate are
\begin{center}
\begin{adjustbox}{max width=\columnwidth}
\(\displaystyle
\begin{aligned}
\admissibleHistorical(\CloseCut(e),h_C,C)\iff{}&
 h_C=\hash(\mathsf{RAI/CloseCut}\parallel e\parallel\operatorname{canon}(C))\\
&\wedge\func{CutSupported}(e,h_C),\\
\admissible_{P_j}(\CloseCut(e),h_C,C)\iff{}&
 \admissibleFresh_{P_j}(\CloseCut(e),h_C,C)\\
&\vee\admissibleHistorical(\CloseCut(e),h_C,C).
\end{aligned}
\)
\end{adjustbox}
\end{center}
The proof \(\rho\), each start witness, and the signer's own report lock are selected locally from the
report and vote store.  They are not serialized with \(C\), and close votes continue to sign only the
close hash rather than a witness bundle reference.  Additional non-conflicting reports may change the
replica's current visible set, and newly learned report equivocation removes that signer's weight from
fresh report calculations.  A retained version is freshly revalidated before unsupported voting; once
\(\func{CutSupported}(e,h_C)\) holds, the historical branch is permanent and later report-store changes
cannot invalidate the cut.  Historical cut support may override a signer's fresh cut lock, but it does
not release the reported election: any unresolved report/finality obligation is carried into fresh
close-record validation.

For a cut election \(q\), let \(\Sigma_{j,q}\) be all valid logical vote leaves for \(q\) currently known to
\(P_j\).  Fresh close-record construction applies the deterministic resolver
\[
\operatorname{resolve}(q,\Sigma_{j,q})=
\begin{cases}
\FinalizedState(b,\pi),
  & \Sigma_{j,q}\text{ derives valid global fast or final evidence for }b,\\
\ReleasedState(\pi_{\mathsf{timeout}}),
  & \Sigma_{j,q}\text{ derives valid timeout evidence},\\
\ReleasedState(\pi_{\mathsf{conflict}}),
  & \Sigma_{j,q}\text{ derives conflicting notarization evidence},\\
\SelectedState(q,b,\pi_{\mathsf{notar}}),
  & b\text{ is the unique known non-timeout notarized value}\\
  & \text{and no final or release evidence is currently known},\\
\bot, & \text{otherwise.}
\end{cases}
\]
The symbols \(\pi\) denote locally selected authenticated vote-leaf witnesses, not certificate objects.

Once the decided cut is installed, it fixes the immutable drain input
\begin{center}
\begin{adjustbox}{max width=\columnwidth}
\(\displaystyle
\mathsf{CloseInput}_e=
\bigl(\closeHash[e-1],\{\frontier[e-1,a]:a\in\mathsf{Accounts}\},\closeCut[e]\bigr).
\)
\end{adjustbox}
\end{center}
This is an objective protocol boundary rather than a replica-local wall-clock snapshot.  Starting from
the preceding decided frontiers, the resolver is applied only to elections in \(\closeCut[e]\).  In
addition, every currently known ordinary-finalized election \(q=\Slot(a,s,e)\), including one outside
the cut, contributes its exact certified segment.  Elections of the concurrently open epoch \(e+1\)
never contribute to epoch \(e\)'s frontier.

For a frontier map \(\Phi\), define \(\LedgerIncludesMap(\Phi,q,b)\) exactly as
\(\func{LedgerIncludes}_e(q,b)\), but with membership tested against the account chains ending at
\(\Phi\).  The fresh close-record lock of correct replica \(P_j\) is
\[
\begin{aligned}
\recordSafe_j(e,\Phi)\iff{}&
\forall q,b:\ordinaryFinalized_j(q,b)\Rightarrow\LedgerIncludesMap(\Phi,q,b)\\
&\wedge\forall q=\Slot(a,s,e):q\in R_j[e]\Rightarrow\\
&\qquad\Bigl(\cannotFinalize_j(q)\vee\exists b\neq\timeoutVal_q:
\LedgerIncludesMap(\Phi,q,b)\\
&\qquad\qquad\wedge\forall Q\in\SignerCommittees(P_j,q):
\func{Support}(P_j,Q,q)\subseteq\{b\}\Bigr).
\end{aligned}
\]
Thus a signer that already ordinary-finalized \(b\) may omit \(q\) from its report but never freshly
votes for a close record omitting \(b\).  A signer that reported \(q\) freshly requires either a
support-compatible included segment or positive evidence that no ordinary finality can exist.

Define persistent close-record support by
\begin{center}
\begin{adjustbox}{max width=\columnwidth}
\(\displaystyle
\begin{aligned}
\RecordSupported(e,h_F)\iff{}&\exists r:\ \forall Q\in\Committees(\CloseRecord(e,r)),\\
&\exists k_Q\in\{\notarKind,\fastKind,\finalKind\}:\
\Cert_Q(k_Q,\CloseRecord(e,r),h_F).
\end{aligned}
\)
\end{adjustbox}
\end{center}
A reconstructed frontier map \(\Phi\) is admissible for \(h_F\) at correct replica \(P_j\) when
\[
\begin{aligned}
\admissibleFresh_{P_j}(\CloseRecord(e),h_F,\Phi)\iff{}&
 h_F=\hash(\mathsf{RAI/CloseRecord}\parallel e\parallel\closeHash[e-1]
       \parallel\operatorname{canon}(\Phi))\\
&\wedge\func{validatedFrontiers}(e,\closeCut[e],\Phi)
\wedge\recordSafe_j(e,\Phi),\\
\admissibleHistorical(\CloseRecord(e),h_F,\Phi)\iff{}&
 h_F=\hash(\mathsf{RAI/CloseRecord}\parallel e\parallel\closeHash[e-1]
       \parallel\operatorname{canon}(\Phi))\\
&\wedge\RecordSupported(e,h_F)
\wedge\func{validatedFrontiers}(e,\closeCut[e],\Phi),\\
\admissible_{P_j}(\CloseRecord(e),h_F,\Phi)\iff{}&
\admissibleFresh_{P_j}(\CloseRecord(e),h_F,\Phi)\\
&\vee\admissibleHistorical(\CloseRecord(e),h_F,\Phi).
\end{aligned}
\]
The predicate \(\func{validatedFrontiers}\) obtains every block on every account chain from genesis
through the claimed frontiers, validates every block and transition, preserves every block committed by
the preceding decided frontiers, and permits at most one block per logical account slot.  For every
\(q\in\closeCut[e]\), a segment tip included on \(\Phi\) must have locally derivable global fast/final
evidence or a non-timeout notarization certificate for that same segment tip; an omitted cut election
must have locally derivable timeout or conflicting terminal-certificate evidence.  A segment included
for \(q=\Slot(a,s,e)\notin\closeCut[e]\) must have locally derivable global fast or final evidence for
that exact tip.  Every strict ancestor of a selected or ordinary-finalized segment tip beyond the
preceding certified frontier must occur in that candidate's \(\SegmentBlocks(q,b)\) suffix or in an
earlier epoch-\(e\) selected or ordinary-finalized segment on the same account chain.  Only epoch-\(e\)
elections may advance an account beyond \(\frontier[e-1,\cdot]\); ordinary finalizations from epoch
\(e+1\) first become eligible for the next close record.  Additional evidence may change a fresh
canonical frontier map.  Once \(\RecordSupported(e,h_F)\) holds, the historical branch is permanent:
by the lock-intersection lemma below, no omitted ordinary-finalized epoch-\(e\) segment can later be
revealed consistently with that support.  Admissibility permits voting only and does not finalize or
release anything before the close-record decision.  In algorithms executed at replica \(P_j\), the
unsubscripted notation \(\admissible\) abbreviates the corresponding replica-local
\(\admissible_{P_j}\) predicate.

\paragraph{Hash reconciliation of close versions.}
For \(T\in\{\mathsf{cut},\mathsf{frontiers}\}\), a replica retains canonical preimages indexed by
hash.  A requester that possesses base hash \(a\) and needs target hash \(b\) sends
\[
\ReconReq(T,e,a,b).
\]
A responder may answer only if it possesses both canonical preimages.  For cuts it sends
\[
\ReconDelta(\mathsf{cut},e,a,b,\Delta^+,\Delta^-),
\qquad
\Delta^+=C_b\setminus C_a,
\quad
\Delta^-=C_a\setminus C_b.
\]
The requester reconstructs
\[
C'=(C_a\setminus\Delta^-)\cup\Delta^+,
\]
checks the canonical target hash, validates \(\admissible(\CloseCut(e),b,C')\), and then caches
\(\cutVersion[e,b]=C'\).  Reconstructing a historical cut does not alter its current report store, vote store, or derived visibility.

For frontier maps the responder sends sparse per-account replacements
\[
\ReconDelta(\mathsf{frontiers},e,a,b,
\{(x,f_a(x),f_b(x)):f_a(x)\neq f_b(x)\}).
\]
The requester checks every expected old frontier, replaces only the named account entries, obtains any
missing blocks through ordinary block gossip, validates the affected chains and slot evidence, recomputes
\[
b=\hash(\mathsf{RAI/CloseRecord}\parallel e\parallel\closeHash[e-1]
\parallel\operatorname{canon}(\Phi_b)),
\]
and caches \(\frontierVersion[e,b]=\Phi_b\).  A replacement may move a pre-decision candidate frontier
forward, backward, or sideways; it never deletes a block from the immutable block store.  A responder
that lacks the base or target sends \(\ReconMiss(T,e,a,b)\), after which the requester asks another
replica or obtains a complete canonical version.  Reconciliation messages are unauthoritative transport
metadata: acceptance always requires recomputing the target hash and applicable admissibility validation.

A correct replica retains every close version for which it votes until the logical close instance
decides.  Decided cut and frontier versions, or delta paths from retained checkpoints, remain available
from correct archival replicas so that joining replicas can reconstruct the complete close history.

There is no distinguished validator proposer and no explicit proposal object in the election core.
For a slot election, the account owner creates and authorizes a contiguous candidate segment ending at \(B_t\);
the owner or any relay may gossip its blocks, and a replica computes \(b=\hash(B_t)\) only after
validating every block and the certified base.  Validators do not synthesize account blocks.

A correct replica constructs its fresh close-cut version from its current report-and-vote visibility set after
locally deriving a valid report quorum.  Its fresh close-cut hash is
\[
\FreshCloseHash_{P_j}(\CloseCut(e))
=
\hash(\mathsf{RAI/CloseCut}\parallel e\parallel\operatorname{canon}(C_{e,j})),
\]
and it retains the canonical set under that hash.

After the decided cut is drained, a correct replica first derives every global fast or final segment-tip
certificate available for elections \(q=\Slot(a,s,e)\), whether or not \(q\in\closeCut[e]\), and runs
\(\FinalizeSegment(q,b)\) for each such certified segment in canonical account-chain order.  It then
constructs a fresh close-record version from \(\mathsf{CloseInput}_e\): a close-local replay starts from
\(\closeState[e-1]\), applies the resolver only to elections in the decided cut, overlays every locally
ordinary-finalized epoch-\(e\) segment outside the cut, and performs validated ancestry closure.  It does
not read ordinary finalizations from epoch \(e+1\).  This derives a complete canonical frontier map
\(\Phi_{e,j}\) satisfying \(\recordSafe_j(e,\Phi_{e,j})\).  Define
\[
\FreshCloseHash_{P_j}(\CloseRecord(e))
=
\hash(\mathsf{RAI/CloseRecord}\parallel e\parallel\closeHash[e-1]
\parallel\operatorname{canon}(\Phi_{e,j})),
\]
and retain the canonical frontier map under that hash.
Fresh construction is used only when the immediately preceding close round is positively proved dead
or when starting round \(0\).  A live carried hash overrides fresh preference.

For \(r>0\), correct replica \(P_j\) chooses the first-vote value for \(I_r\) as follows:
\[
\func{NextCloseHash}_{P_j}(I,r)=
\begin{cases}
\FreshCloseHash_{P_j}(I),
  & \CloseDead_{P_j}(I_{r-1}),\\
h,
  & \LiveCarry_{P_j}(I_{r-1},h),\\
\bot,
  & \text{otherwise.}
\end{cases}
\]
For round \(0\),
\[
\func{NextCloseHash}_{P_j}(I,0)=
\FreshCloseHash_{P_j}(I).
\]
Death evidence takes precedence if the replica knows both a non-timeout notarization certificate and
a proof that the source round is dead.  If the result is \(\bot\), the replica does not start the
successor round.

A correct replica never abandons a carried hash merely because:
\begin{itemize}
\item its fresh close hash changes;
\item it learns additional valid close evidence;
\item another admissible hash becomes freshly preferred;
\item its local timer expires; or
\item it has already entered another local protocol state.
\end{itemize}
A carry may be abandoned only after a valid death proof for the round that created that carry.  Thus
\(\text{carry overrides fresh preference}\).

For a fixed account slot \((a,s)\), retries are sequential.  A replica does not start or first-vote in a later
retry \(\Slot(a,s,e')\) while some earlier election \(\Slot(a,s,e_0)\), \(e_0<e'\), remains pending.
A later same-account-slot retry may start only after the certified close record has been reconstructed and installed in the
normative order below, the resulting close state records \(\ReleasedState(\pi)\) for the earlier
election, and \(\finalizedSlot[a,s]=\bot\).  If the earlier segment is finalized either by an ordinary
certificate or by the close record, no retry of that starting slot is started.  Adjacent retries of the same starting slot remain
constrained by shared-committee vote locks until certified release, and disjoint future elections are
constrained by the governing certified finalized frontier.

Local segment completion is
\begin{center}
\begin{adjustbox}{max width=\columnwidth}
\(\displaystyle
\begin{aligned}
\Complete(\Slot(a,s,e),b)
\iff{}& \NotarCert(\Slot(a,s,e),b)\in\mathsf{pool}
\wedge b\neq\timeoutVal_{\Slot(a,s,e)}\\
&\wedge\SegmentBlocks(\Slot(a,s,e),b)=(B_s,\ldots,B_t)\\
&\wedge\forall r\in[s,t]:\PayloadAvailable(B_r)
\wedge\SegmentValid(\Slot(a,s,e),b).
\end{aligned}
\)
\end{adjustbox}
\end{center}
A completed segment inserts every missing non-genesis block into \(\completeTree\) together:
\begin{center}
\begin{adjustbox}{max width=\columnwidth}
\(\displaystyle
\frac{\Complete(q,b)\quad\SegmentBlocks(q,b)=(B_s,\ldots,B_t)}
{\completeTree\gets\completeTree\cup\{B_r:s\le r\le t\}}.
\)
\end{adjustbox}
\end{center}
Hash deduplication makes the insertion idempotent.  No message is broadcast solely because blocks were
inserted.

Epoch membership is not stored on every block.  Instead, a fresh close-record version derives one
frontier for every account and hashes the canonical frontier map together with the preceding decided
close-record hash:
\begin{center}
\begin{adjustbox}{max width=\columnwidth}
\(\displaystyle
\begin{aligned}
\Phi_{e,j}
    &=\operatorname{sort}\{(a,f_{e,j}(a)):a\in\mathsf{Accounts}\},\\
R^{\mathsf{ledger}}_{e,j}
    &=\hash(\operatorname{canon}(\Phi_{e,j})),\\
h_{F,j}
    &=\hash(\mathsf{RAI/CloseRecord}\parallel e\parallel\closeHash[e-1]
       \parallel\operatorname{canon}(\Phi_{e,j})).
\end{aligned}
\)
\end{adjustbox}
\end{center}
The decided close cut is supporting validation context for deriving \(\Phi_{e,j}\), but it is not an
additional field in the close-record value and is not itself the durable ledger-consensus object.  The
immutable \(\mathsf{CloseInput}_e\) consists of the preceding certified close, the preceding decided
frontiers, and the decided drain cut; ordinary-finalized epoch-\(e\) evidence remains a monotone input
outside that cut.  A close-record version is replica-relative before the logical close instance decides;
once \(\closeDecision[\CloseRecord(e)]\) is set, write the decided frontier map without the replica
subscript as \(\Phi_e\).  Blocks, reports, and vote batches are obtained through ordinary gossip, while
the frontier map itself is reconstructed by hash reconciliation when necessary.

Because every block contains its predecessor hash, each frontier transitively commits to its entire
account chain.  The validator recomputes every block hash and owner signature, checks account and
sequence continuity, replays every debit and credit, verifies finalized send-before-receive dependencies,
checks destination ownership and unique send consumption, applies every representative change, and
rejects any chain with a missing or invalid predecessor.  The replay begins from the preceding certified close state and applies blocks selected from the decided
epoch-\(e\) cut together with every validated ordinary-finalized epoch-\(e\) segment, including segments
whose elections are outside the cut.  Ordinary finalizations from epoch \(e+1\) do not affect this replay.  Every receive
must resolve to a still-pending finalized send addressed to the receiving account.  It also requires
every block on the preceding certified account chains to remain, every balance to remain non-negative,
and the supply invariant to hold.
The resulting \(\bal_e\), \(\repOf_e\), \(\pendingSend_e\), and \(\receivedSend_e\) maps are part of
the reconstructed \(\closeState[e]\).  The canonical hash of the complete frontier map is
\(R^{\mathsf{ledger}}_e\), while the close-record hash additionally commits to epoch \(e\) and
\(\closeHash[e-1]\).

Let \(\frontier[e,a]\) be the decided frontier of account \(a\) through epoch \(e\), with
\(\frontier[e,a]=\frontier[e-1,a]\) for accounts that do not move in epoch \(e\).  The blocks of
account \(a\) belonging to epoch \(e\) are the ordered account-chain segment after
\(\frontier[e-1,a]\) and up to \(\frontier[e,a]\):
\begin{center}
\begin{adjustbox}{max width=\columnwidth}
\(\displaystyle
\EpochBlocks(e,a)=
(\frontier[e-1,a],\frontier[e,a]] .
\)
\end{adjustbox}
\end{center}
Only the frontier block is named in \(\Phi_e\); every intermediate block is committed through the
predecessor-hash chain and its epoch membership is obtained from account-tree order.

For \(q=\Slot(a,s,e)\), the decided ledger includes candidate tip \(b\) exactly when
\(\func{LedgerIncludes}_e(q,b)\).  Every block in \(\SegmentBlocks(q,b)\), together with every earlier
ancestor on its certified account chain, is finalized.  Every other candidate of \(q\) is released.  If no candidate of \(q\) occurs on the
certified ledger, then \(q\) is released.  An epoch-\(e\) slot election outside the certified cut is
finalized when its ordinary-certified segment occurs on the decided ledger and is otherwise released by
certified close-record omission.  Before decision, ordinary-finalized and selected are evidence
classifications only; after certification the durable distinction is whether a block occurs on a chain
committed by the decided frontier map.

Once the logical close-record instance decides \(h_F\), a replica installs the corresponding
reconstructed frontier version atomically in the following normative order: (1) locally derive the
deciding fast or final certificate from signed close-record vote leaves; (2) reconstruct \(\Phi_e\) for \(h_F\) by a
retained version, reconciliation delta, or complete version transfer; (3) verify
\[
h_F=\hash(\mathsf{RAI/CloseRecord}\parallel e\parallel\closeHash[e-1]
\parallel\operatorname{canon}(\Phi_e));
\]
(4) reconstruct and historically validate the decided close-cut version and locally derive every required
cut-election certificate, timeout, or conflict, together with every ordinary fast/final certificate used
for an epoch-\(e\) segment outside the cut, from signed vote leaves; (5) obtain and validate every block on every account
chain from genesis through the claimed frontiers; (6) verify the frontier classifications, ancestry
closure, ledger transitions, and \(R^{\mathsf{ledger}}_e=\hash(\operatorname{canon}(\Phi_e))\); (7)
run \(\FinalizeSegment(q,b)\) for every newly included segment obligation in canonical account-chain order; and only then (8)
derive and install \(\closeState[e]\), \(\frontier[e,\cdot]\),
\(\ledgerRoot[e]\gets R^{\mathsf{ledger}}_e\), and \(\closeHash[e]\gets h_F\).  No retry or
future-vote decision may read the new release state before all eight steps complete.

A joining replica may wait for the next live close to learn the latest synchronization target, but that
observation is not a trust assumption.  It starts from genesis, obtains the durable blocks, reports, and
vote batches, and replays these steps for every epoch in order.  It reconstructs each decided cut and frontier
preimage, validates every block and state transition, locally derives each deciding vote quorum,
recomputes every close-record hash from \(\closeHash[e-1]\) and \(\Phi_e\), and accepts epoch \(e\)
only when the recomputed hash is the value decided by the close-record votes.  The latest validated
close-record hash is the current consensus checkpoint: its predecessor-hash chain and account
frontiers transitively commit to every block on the certified ledger chains.  Thus the replica derives
the live committee internally and bootstraps the entire ledger without trusting an archive's derived
state or an externally supplied committee snapshot.

If \(\FirstCert(q,b)\in\mathsf{pool}\), \(\FinalCert(q,b)\in\mathsf{pool}\), or a decided
ledger-state root includes \(b\) for \(q=\Slot(a,s,e)\), with
\(b\neq\timeoutVal_q\) and \(\Complete(q,b)\), then \(\FinalizeSegment(q,b)\) runs.  It requires
\(\StablePrefix(q,B_t)\wedge\SegmentValid(q,b)\) and reconstructs
\(\SegmentBlocks(q,b)=(B_s,\ldots,B_t)\).  The replica copies the finalized ledger and affected
\(\finalizedSlot\) entries into temporary state, then validates and applies every block in
ancestor-to-descendant order.  For every \(r\in[s,t]\), it asserts
\(\finalizedSlot[a,r]\in\{\bot,\hash(B_r)\}\); a different value is a safety fault.  No globally visible
ledger, frontier, release, or finalized-slot write occurs during this replay.  Only after every block,
receive, balance, representative update, and the supply invariant succeeds does one atomic commit install
all temporary ledger changes and set
\[
\forall r\in[s,t]:\quad\finalizedSlot[a,r]\gets\hash(B_r).
\]
A crash or validation failure before the commit leaves the complete pre-segment state visible.  Timeout
and release outcomes never invoke \(\FinalizeSegment\).

Normative invariants: enabled voting; committee-local first-vote uniqueness; certified release for conflicting slot
vote history; timeout locality; separation of notarization from slot finality; no gaps in
\(\completeTree\); at most one complete-tree entry per block hash; certified per-account epoch
frontiers; a canonical frontier map and hash-chained close-record value over every account frontier; complete validation of every
account chain from genesis through its frontier; valid owner signatures; non-negative balances; finalized
send-before-receive dependencies; destination-only and at-most-once send consumption; deterministic
representative updates; supply conservation; monotonic \(\finalizedSlot\); atomic segment finality by ordinary fast or final certificates or by certified ledger inclusion from a certified base; release of every epoch election omitted from the certified ledger; inclusion of every ordinary-finalized epoch-\(e\) segment in any freshly admissible close record; fresh report locks for close-cut and close-record voting; persistent historical admissibility only after same-hash certificate support; close-cut boundedness after a
certified close-cut hash; a locally derivable single-committee report quorum and valid reconstructed cut for every non-timeout close-cut vote;
mandatory validation under \(\closeState[e-2]\) for every epoch-\(e\) vote; deterministic stake-weighted committee
derivation from certified balances, representatives, close state, and close-record hash; uniqueness and irrevocability of the certified close
record; and at most one open epoch and at most one closing epoch, with older epochs closed.

\section{Visibility and Close Cuts}

\begin{center}
\begin{adjustbox}{max width=\columnwidth}
\(\displaystyle
\begin{aligned}
\Reportable_j(\Slot(a,s,e))\iff{}& \started[\Slot(a,s,e)]\wedge
\neg\cannotFinalize_j(\Slot(a,s,e))\\
&\wedge\neg\exists b:\ordinaryFinalized_j(\Slot(a,s,e),b)\\
&\wedge\Bigl(\neg\finished(\Slot(a,s,e))\vee
\exists Q\in\Committees(\Slot(a,s,e)),\ \exists b\neq\timeoutVal_{\Slot(a,s,e)}:\\
&\qquad b\in\func{Support}(P_j,Q,\Slot(a,s,e))\Bigr).
\end{aligned}
\)
\end{adjustbox}
\end{center}
A correct report therefore retains every started election that may still ordinary-finalize and every
supporting non-timeout obligation not already discharged by positive non-finalizability evidence.  An
election already ordinary-finalized at the reporter is settled and need not appear in that report.  The
omission does not weaken close-record safety because the reporter's exact ordinary-finality lock replaces
the report lock.

For every signer and epoch, \(\reportStore[P_i,e]\) retains every distinct valid signed report learned
for that signer.  Report equivocation and eligibility are derived objectively as
\[
\begin{aligned}
\ReportEquivocated(P_i,e)\iff{}&\exists R\neq R':
\Report(e,R,\sigma),\Report(e,R',\sigma')\in\reportStore[P_i,e],\\
\EligibleReport(P_i,e,R)\iff{}&\Report(e,R,\sigma)\in\reportStore[P_i,e]
\wedge\neg\ReportEquivocated(P_i,e).
\end{aligned}
\]
The derived value \(report[P_i,e]\) is the unique body \(R\) satisfying
\(\EligibleReport(P_i,e,R)\), and is \(\bot\) otherwise.  Thus all signed reports are retained,
but an equivocating signer contributes no weight to report quorums or report-derived visibility.

A report proof \(\rho\) contains distinct eligible signed reports from the deterministic reporting
committee \(Q^\star=\ReportCommittee(e)\), with combined signer weight crossing the close barrier:
\begin{center}
\begin{adjustbox}{max width=\columnwidth}
\(\displaystyle
\ReportProof(e,\rho)\iff
Q^\star=\ReportCommittee(e)\neq\bot\wedge
\SignerWeight_{Q^\star}
\bigl(\{P_i\in Q^\star:\exists R_i,\ \Report(e,R_i,\sigma_i)\in\rho
\wedge\EligibleReport(P_i,e,R_i)\}\bigr)\ge W_{Q^\star}-F_{Q^\star}.
\)
\end{adjustbox}
\end{center}
All reports in \(\rho\) must be individually valid and eligible.  An equivocating signer contributes no
weight.  The account slots forced by the proof are
\begin{center}
\begin{adjustbox}{max width=\columnwidth}
\(\displaystyle
\ReportVisible(e,\rho)=
\{(a,s):
\SignerWeight_{Q^\star}
(\{P_i\in Q^\star:\exists R_i,\ \Report(e,R_i,\sigma_i)\in\rho
\wedge\EligibleReport(P_i,e,R_i)\wedge\Slot(a,s,e)\in R_i\})>F_{Q^\star}\}.
\)
\end{adjustbox}
\end{center}
Because every global non-timeout finishing certificate still contains a compatible local certificate
from \(Q^\star\), every correct certificate signer either reports the unfinished obligation or already
holds an exact ordinary-finality lock.  Report-derived visibility continues to force sufficiently widely
reported unresolved elections into the cut, while the combined report/finality locks protect ordinary
finality at the close-record phase.  Slot and close finality continue to require every election committee.

\begin{breakablealgorithm}
\caption{\(\UpdateVisible(e)\)}
\begin{algorithmic}[1]
\State $Q^\star\gets\ReportCommittee(e)$; $V_{\mathrm{reports}}\gets\{(a,s):\SignerWeight_{Q^\star}(\{P_i\in Q^\star:\Slot(a,s,e)\in report[P_i,e]\})>F_{Q^\star}\}$
\State $V_{\mathrm{votes}}\gets\{(a,s):\exists Q\in\Committees(\Slot(a,s,e)):\mathsf{pool}\text{ contains valid votes of distinct-signer weight }>F_Q\text{ in }Q\text{ for }\Slot(a,s,e)\}$
\State $V\gets V_{\mathrm{reports}}\cup V_{\mathrm{votes}}$
\If{$V\neq visible[e]$}
    \State $visible[e]\gets V$
    \If{$state[e]=\mathsf{closing}$ and a cached valid report proof $\rho$ exists}
        \State $C\gets\{\Slot(a,s,e):(a,s)\in visible[e]\cup\ReportVisible(e,\rho)\}\cup\{q:q=\Slot(a,s,e)\wedge\reportLocked_j(q)\}$
        \State $h_C\gets\hash(\mathsf{RAI/CloseCut}\parallel e\parallel\operatorname{canon}(C))$
        \State $\cutVersion[e,h_C]\gets C$
    \EndIf
\EndIf
\end{algorithmic}
\end{breakablealgorithm}

Reports contain pending or otherwise still-finalizable slot elections, but omit elections already
ordinary-finalized by the reporter and elections with positive non-finalizability evidence.  Each
reported election creates a signer-local fresh close lock until non-finalizability is proved; learning
ordinary finality changes it to an exact frontier-inclusion lock.  Report-derived visibility is evaluated
only in the deterministic reporting committee, while
vote-derived visibility is evaluated in any election committee.  Report-derived visibility is recomputed
from currently eligible reports and may shrink when report equivocation is learned; vote-derived visibility
remains monotone because valid votes are persistent.  A replica's current close cut is
\(C=\{\Slot(a,s,e):(a,s)\in visible[e]\}\), and its close-cut hash is
\[
\hash(\mathsf{RAI/CloseCut}\parallel e\parallel\operatorname{canon}(C)),
\]
not a hash of reports and not the durable epoch close-record hash.  A slot election enters current visibility either through report signer weight greater than the Byzantine bound in \(\ReportCommittee(e)\), or through visible valid-vote signer weight greater than \(F_Q\) in some election committee.  Before a non-timeout fresh close-cut vote, the replica locally selects a valid report proof \(\rho\),
requires the canonical cut to contain its current visible set, every account slot in
\(\ReportVisible(e,\rho)\), and every election under its own fresh report lock, and validates a local
\(\func{StartWitness}_e(q)\) for every
\(q\in C\), and stores \(\cutVersion[e,h_C]=C\).  No cut, proof, witness map, or certificate is
broadcast; only the close vote containing \(h_C\) is disseminated.

If \(\closeDecision[\CloseCut(e)]=(r,h_C)\), the replica reconstructs and historically validates
\[
C=\cutVersion[e,h_C]
\]
by retained state or hash reconciliation and sets \(\closeCut[e]=C\).  This fixes
\(\mathsf{CloseInput}_e=(\closeHash[e-1],\frontier[e-1,\cdot],C)\).  Its current visibility may differ
from \(C\); reconstructing the decided historical version does not remove reports or votes and does not
overwrite the visibility derived from those stores.  For every \(\Slot(a,s,e)\in C\), each replica resolves the frozen drain obligation.  It derives a fresh
close-record frontier map by close-local replay from the preceding certified state using that cut
evidence together with every currently known ordinary-finalized epoch-\(e\) segment outside \(C\) and
the validated block tree.
Different correct replicas may initially derive different admissible hashes because their evidence for
the frozen elections differs.  Newly gossiped atoms may change fresh preference but cannot invalidate a
previously validated historical version or override a live carry.  Replicas vote on
\[
h_F=\hash(\mathsf{RAI/CloseRecord}\parallel e\parallel\closeHash[e-1]
\parallel\operatorname{canon}(\Phi_e)).
\]
Only the value stored in \(\closeDecision[\CloseRecord(e)]\) becomes \(\closeHash[e]\); the close-cut
hash and local slot-drain outcomes do not.

\section{Main Algorithms}

\begin{breakablealgorithm}
\caption{RAI initialization for replica $P_j$}
\begin{algorithmic}[1]
\State $e_{\mathrm{open}}\gets0$; $state[-1]\gets\mathsf{closed}$; $state[e_{\mathrm{open}}]\gets\mathsf{open}$; let $G$ be the genesis epoch state
\State $\closeState[-1]\gets G$; $\closeHash[-1]\gets\hash(G)$; $\committee_{\mathsf{gen}}\gets\deriveCommittee(\mathsf{genesis},G)$; $\committee[\cdot]\gets\bot$
\State $\completeTree\gets\{\mathsf{genesis}_a:a\in\mathsf{Accounts}\}$; $\frontier[-1,a]\gets\mathsf{genesis}_a$ for every account $a$; derive $\bal_{-1}$, $\repOf_{-1}$, $\pendingSend_{-1}$, and $\receivedSend_{-1}$ from $G$
\State $\ledgerRoot[-1]\gets\hash(\operatorname{sort}\{(a,\frontier[-1,a]):a\in\mathsf{Accounts}\})$
\State $\finalizedSlot[\cdot,\cdot]\gets\bot$
\State $epochStart[\cdot]\gets\bot$; $epochStart[e_{\mathrm{open}}]\gets\func{clock}()$; $T_{\mathrm{election}}[\cdot]\gets\bot$
\State $outcome[\cdot]\gets\pending$; $\closeDecision[\cdot]\gets\bot$; $\closeRound[\cdot]\gets0$; $notarized[\cdot,\cdot,\cdot]\gets\{\}$; $secondLook[\cdot,\cdot,\cdot]\gets\{\}$
\State $firstVote[\cdot,\cdot,\cdot]\gets\bot$; $finalVoted[\cdot,\cdot,\cdot]\gets\false$; $\votedFor[\cdot,\cdot,\cdot]\gets\false$; $\blockMap[\cdot,\cdot]\gets\bot$
\State $visible[e_{\mathrm{open}}]\gets\{\}$; $\reportStore[\cdot,e_{\mathrm{open}}]\gets\{\}$; $\closeHash[e_{\mathrm{open}}]\gets\bot$; $\ledgerRoot[e_{\mathrm{open}}]\gets\bot$; $\closeState[e_{\mathrm{open}}]\gets\bot$; $\cutVersion[e_{\mathrm{open}}]\gets\{\}$; $\frontierVersion[e_{\mathrm{open}}]\gets\{\}$; $\closeCut[e_{\mathrm{open}}]\gets\bot$
\end{algorithmic}
\end{breakablealgorithm}

\begin{breakablealgorithm}
\caption{\(\StartClosing(k)\)}
\begin{algorithmic}[1]
\State \textbf{require} $state[k]=\mathsf{open}$ and $state[k-1]=\mathsf{closed}$ and $\closeState[k-1]\neq\bot$
\State $state[k]\gets\mathsf{closing}$
\Statex \Comment{The epoch-\(k\) drain input is fixed by the decided cut and \(\frontier[k-1,\cdot]\); ordinary-finalized epoch-\(k\) segments outside the cut may still enter the close record, but epoch \(k+1\) segments may not.}
\State $\closeDecision[\CloseCut(k)]\gets\bot$; $\closeRound[\CloseCut(k)]\gets0$; $\closeDecision[\CloseRecord(k)]\gets\bot$; $\closeRound[\CloseRecord(k)]\gets0$
\If{$P_j\in\ReportCommittee(k)$}
    \State $R_j[k]\gets\{\Slot(a,s',k):a\in\mathsf{Accounts},\ s'\ge1,\ \Reportable_j(\Slot(a,s',k))\}$
    \State $\reportStore[P_j,k]\gets\reportStore[P_j,k]\cup\{\Report(k,R_j[k],\sigma_j)\}$; broadcast $\Report(k,R_j[k],\sigma_j)$
\EndIf
\State $e_{\mathrm{open}}\gets k+1$; $state[e_{\mathrm{open}}]\gets\mathsf{open}$; $epochStart[e_{\mathrm{open}}]\gets\func{clock}()$
\State $visible[e_{\mathrm{open}}]\gets\{\}$; $\reportStore[\cdot,e_{\mathrm{open}}]\gets\{\}$; $\closeHash[e_{\mathrm{open}}]\gets\bot$; $\ledgerRoot[e_{\mathrm{open}}]\gets\bot$; $\closeState[e_{\mathrm{open}}]\gets\bot$; $\cutVersion[e_{\mathrm{open}}]\gets\{\}$; $\frontierVersion[e_{\mathrm{open}}]\gets\{\}$; $\closeCut[e_{\mathrm{open}}]\gets\bot$
\State load the certified finalized frontiers and release evidence from $\closeState[k-1]$ as durable validation state
\end{algorithmic}
\end{breakablealgorithm}

\begin{breakablealgorithm}
\caption{\(\FinishClosing(k)\)}
\begin{algorithmic}[1]
\State \textbf{require} $state[k]=\mathsf{closing}$ and $state[k-1]=\mathsf{closed}$
\State $state[k]\gets\mathsf{closed}$; $\func{clearEpochVars}(k)$
\end{algorithmic}
\end{breakablealgorithm}

\begin{breakablealgorithm}
\caption{\(\func{ObtainCloseVersion}_{P_j}(I,h)\)}
\begin{algorithmic}[1]
\If{$I=\CloseCut(e)$ and $\cutVersion[e,h]=C$ and $\admissible(I,h,C)$}
    \State \Return $C$
\ElsIf{$I=\CloseRecord(e)$ and $\frontierVersion[e,h]=\Phi$ and $\admissible(I,h,\Phi)$}
    \State \Return $\Phi$
\EndIf
\ForAll{locally known base hashes $a$ for the same close instance}
    \State send $\ReconReq(T,e,a,h)$, where $T=\mathsf{cut}$ or $\mathsf{frontiers}$
    \If{a valid $\ReconDelta(T,e,a,h,\delta)$ is received}
        \State apply $\delta$ to the retained preimage of $a$
        \State recompute exactly $h$ and run the applicable admissibility checks
        \If{validation succeeds}
            \State cache and \Return the reconstructed version
        \EndIf
    \EndIf
\EndFor
\State retrieve a complete canonical version for $h$
\State recompute exactly $h$ and run the applicable admissibility checks
\State cache and \Return the version
\end{algorithmic}
\end{breakablealgorithm}

\begin{breakablealgorithm}
\caption{Start close round \(I_r\) at correct replica \(P_j\)}
\begin{algorithmic}[1]
\State \textbf{require} $\closeDecision[I]=\bot$
\If{$r=0$}
    \State $x\gets\FreshCloseHash_{P_j}(I)$
\ElsIf{$\CloseDead_{P_j}(I_{r-1})$}
    \State $x\gets\FreshCloseHash_{P_j}(I)$
\ElsIf{$\exists h:\LiveCarry_{P_j}(I_{r-1},h)$}
    \State let $h$ be the locally live carried hash
    \State $S\gets\func{ObtainCloseVersion}_{P_j}(I,h)$
    \State $x\gets h$
\Else
    \State \Return without starting round $r$
\EndIf
\State initialize the committee-local election-kernel state for $I_r$
\State $\closeRound[I]\gets\max(\closeRound[I],r)$
\State first-vote $x$ with omitted scope, which targets every applicable close committee
\end{algorithmic}
\end{breakablealgorithm}

When \(x\) is a carried hash, the replica waits to reconstruct and validate its canonical version rather than
replacing \(x\) with a timeout first vote.  The replica does not issue a timeout first vote in \(I_r\)
merely because the carried version is still being reconstructed; equivalently,
\(\text{no timeout first vote while the required carried version is being reconstructed}\).

After the carried first vote has been issued, the committee-local second-look, notarization,
timeout-notarization, fast, and final rules remain enabled.

Here \(\func{clearEpochVars}(e)\) deletes only volatile caches.  It never deletes durable safety
evidence, including \(\closeHash[e]\), \(\ledgerRoot[e]\), \(\closeState[e]\), \(\closeCut[e]\),
\(\committee[e]\), decided \(\cutVersion\) and \(\frontierVersion\) entries, timeout, conflict,
close-exclusion, or certified-ledger evidence, local unreleased votes, \(\completeTree\),
\(\frontier[\cdot,\cdot]\), or \(\finalizedSlot\).

\begin{breakablealgorithm}
\caption{RAI epoch loop for replica $P_j$}
\begin{algorithmic}[1]
\State \textbf{while} true \textbf{wait until either}
\Statex
\State $state[e_{\mathrm{open}}]=\mathsf{open}$ and $state[e_{\mathrm{open}}-1]=\mathsf{closed}$ and $\func{clock}()>epochStart[e_{\mathrm{open}}]+\Delta_{\mathrm{timeout}}$ $\Rightarrow$
\State \quad $k\gets e_{\mathrm{open}}$; $\StartClosing(k)$
\Statex
\State received valid $\Report(k,R_i[k],\sigma_i)$ from $P_i$ with $\Report(k,R_i[k],\sigma_i)\notin\reportStore[P_i,k]$ $\Rightarrow$
\State \quad $\reportStore[P_i,k]\gets\reportStore[P_i,k]\cup\{\Report(k,R_i[k],\sigma_i)\}$; recompute the derived $report[P_i,k]$; $\UpdateVisible(k)$
\Statex
\State $state[k]=\mathsf{closing}\wedge Q^\star=\ReportCommittee(k)\wedge
\SignerWeight_{Q^\star}(\{P_i\in Q^\star:report[P_i,k]\neq\bot\})\ge W_{Q^\star}-F_{Q^\star}\wedge
\closeDecision[\CloseCut(k)]=\bot\wedge\closeRound[\CloseCut(k)]=0\Rightarrow$
\State \quad assemble a valid report proof $\rho_k$ from the cached reports; $\UpdateVisible(k)$
\State \quad $C\gets\{\Slot(a,s,k):(a,s)\in visible[k]\cup\ReportVisible(k,\rho_k)\}\cup\{q:q=\Slot(a,s,k)\wedge\reportLocked_j(q)\}$
\State \quad $h_C\gets\hash(\mathsf{RAI/CloseCut}\parallel k\parallel\operatorname{canon}(C))$; $\cutVersion[k,h_C]\gets C$
\State \quad start close round $\CloseCut(k,0)$ using $\func{NextCloseHash}_{P_j}(\CloseCut(k),0)$
\Statex
\State $\closeDecision[\CloseCut(k)]=(r,h_C)$ and $\closeCut[k]=\bot$ $\Rightarrow$
\State \quad locally derive the deciding $\FirstCert(\CloseCut(k,r),h_C)$ or $\FinalCert(\CloseCut(k,r),h_C)$ from signed vote leaves
\State \quad $C\gets\func{ObtainCloseVersion}_{P_j}(\CloseCut(k),h_C)$; validate $\admissibleHistorical(\CloseCut(k),h_C,C)$
\State \quad $\closeCut[k]\gets C$; fix $\mathsf{CloseInput}_k=(\closeHash[k-1],\frontier[k-1,\cdot],C)$
\State \quad for every $q\in C$: derive its historical start obligation from the certified cut and, if $T_{\mathrm{election}}[q]=\bot$, set $T_{\mathrm{election}}[q]\gets\func{clock}()$
\State \quad \textbf{wait until} $\forall q=\Slot(a,s,k)\in C:\finished(q)$
\State \quad derive every currently available global fast or final certificate for $q=\Slot(a,s,k)$, whether or not $q\in C$, and run $\FinalizeSegment(q,b)$ for each certified segment in canonical account-chain order
\State \quad derive $\Phi_k$ by close-local replay from $\closeState[k-1]$, resolving $q\in C$ and including every locally ordinary-finalized epoch-$k$ segment outside $C$; ignore epoch-$(k+1)$ ordinary finalizations
\State \quad $R^{\mathsf{ledger}}_k\gets\hash(\operatorname{canon}(\Phi_k))$
\State \quad $h_F\gets\hash(\mathsf{RAI/CloseRecord}\parallel k\parallel\closeHash[k-1]\parallel\operatorname{canon}(\Phi_k))$
\State \quad $\frontierVersion[k,h_F]\gets\Phi_k$
\State \quad start close round $\CloseRecord(k,0)$ using $\func{NextCloseHash}_{P_j}(\CloseRecord(k),0)$
\Statex
\State $\closeDecision[\CloseRecord(k)]=(r,h_F)$ and $\closeHash[k]=\bot$ $\Rightarrow$
\State \quad locally derive the deciding $\FirstCert(\CloseRecord(k,r),h_F)$ or $\FinalCert(\CloseRecord(k,r),h_F)$ from signed vote leaves
\State \quad $\Phi_k\gets\func{ObtainCloseVersion}_{P_j}(\CloseRecord(k),h_F)$; validate $\admissible(\CloseRecord(k),h_F,\Phi_k)$
\State \quad $R^{\mathsf{ledger}}_k\gets\hash(\operatorname{canon}(\Phi_k))$; require $h_F=\hash(\mathsf{RAI/CloseRecord}\parallel k\parallel\closeHash[k-1]\parallel\operatorname{canon}(\Phi_k))$
\State \quad obtain and validate every account chain from genesis through the frontiers in $\Phi_k$
\State \quad replay and verify every owner signature, balance, send, receive, and representative transition; atomically install each newly included account-chain suffix and derive released elections by omission
\State \quad install $\closeState[k]$ with $\bal_k$, $\repOf_k$, $\pendingSend_k$, and $\receivedSend_k$, plus $\frontier[k,\cdot]$ and $\ledgerRoot[k]\gets R^{\mathsf{ledger}}_k$; $\closeHash[k]\gets h_F$; $\committee[k]\gets\deriveCommittee(k,\closeHash[k],\closeState[k])$
\State \quad $\FinishClosing(k)$
\end{algorithmic}
\end{breakablealgorithm}

\paragraph{Vote emission and reception.}
Every algorithmic instruction that logically emits a first, notarization, or final vote appends the
corresponding \(\VoteLeaf\) to the signer's current batch for the applicable governing context.  For a
nonempty target set \(S\subseteq\SignerCommittees(P_i,\id)\), notation
\(\FirstVote(\id,x;S)\), \(\NotarVote(\id,x;S)\), or \(\FinalVote(\id,x;S)\) means that the scope
is omitted when \(S=\SignerCommittees(P_i,\id)\) and otherwise encodes the canonical explicit proper
subset \(S\).  The implementation may flush after a size, byte, or latency bound; flushing signs one
\(\VoteBatch\).  Receiving either a complete batch or a leaf with a valid inclusion proof expands each
leaf through its deterministically derived effective scope before certificate counting.  This batching
rule is semantics-preserving and is omitted from repeated pseudocode lines below.

\begin{breakablealgorithm}
\caption{Generic slot-election loop for $\id=\Slot(a,s,e)$ at replica $P_j$}
\begin{algorithmic}[1]
\State \textbf{while} $\neg\finished(\Slot(a,s,e))$ and $state[e]\neq\mathsf{closed}$ \textbf{wait until either}
\Statex
\State $\id=\Slot(a,s,e)$ and $\exists b:\Complete(\Slot(a,s,e),b)$ $\Rightarrow$
\State \quad reconstruct $\SegmentBlocks(\id,b)=(B_s,\ldots,B_t)$ and insert every missing $B_r$ into $\completeTree$
\State \quad $B_{\mathrm p}[\id]\gets b$
\State \quad $S_j\gets\{Q\in\SignerCommittees(P_j,\id):\neg\finalVoted(P_j,\id,Q)\wedge\func{Support}(P_j,Q,\id)\subseteq\{b\}\}$
\If{$\enabled(\id)$ and $S_j\neq\varnothing$ and $\safeToVote_{P_j}(\Slot(a,s,e),b)$}
    \State $\votedFor[P_j,\Slot(a,s,e),b]\gets\true$
    \State for every $Q\in S_j$: $finalVoted[P_j,\id,Q]\gets\true$
    \State logically emit $\FinalVote(\id,b;S_j)$
\EndIf
\Statex
\State the pool contains a certificate for $(t,\id,x)$ with $\Terminal(t,\id,x)$ $\Rightarrow$
\If{$outcome[\id]=\pending$}
    \State $outcome[\id]\gets\certOutcome(t,x)$
\EndIf
\If{$\id=\Slot(a,s,e)$ and $(\FirstCert(\id,x)\in\mathsf{pool}\vee\FinalCert(\id,x)\in\mathsf{pool})$ and $x\neq\timeoutVal_\id$ and $\Complete(\id,x)$}
    \State $\FinalizeSegment(\id,x)$
\EndIf
\Statex
\State locally derive, reconstruct, or learn the required data for value $x$ for $\id$ $\Rightarrow$
\State \quad let $S_j\neq\varnothing$ be the maximal set of unvoted applicable committees in which $x$ is the selected first-vote value
\State \quad \textbf{require} $S_j\subseteq\SignerCommittees(P_j,\id)$ and $\enabled(\id)$ and $\forall Q\in S_j:\neg\firstVoted(P_j,\id,Q)\wedge\neg\finalVoted(P_j,\id,Q)$
\If{$\id=\Slot(a,s,e)$ and $x=b=\hash(B_t)$ for a received candidate segment}
    \State store every validated block $B_r$ in the suffix as $\blockMap[\id,\hash(B_r)]\gets B_r$
\EndIf
\State \quad \textbf{require} $\admissible(\id,x)$
\If{$\id=\Slot(a,s,e)$}
    \State \quad \textbf{require} $\safeToVote_{P_j}(\Slot(a,s,e),x)$
\EndIf
\If{$T_{\mathrm{election}}[\id]=\bot$}
    \State $T_{\mathrm{election}}[\id]\gets\func{clock}()$
\EndIf
\State \quad for every $Q\in S_j$: $firstVote[P_j,\id,Q]\gets x$
\If{$\id=\Slot(a,s,e)$}
    \State \quad $\votedFor[P_j,\Slot(a,s,e),x]\gets\true$
\EndIf
\State \quad logically emit $\FirstVote(\id,x;S_j)$
\Statex
\State $S_j\gets\{Q\in\SignerCommittees(P_j,\id):\neg\firstVoted(P_j,\id,Q)\wedge\neg\finalVoted(P_j,\id,Q)\}$ and $S_j\neq\varnothing$ and $T_{\mathrm{election}}[\id]\neq\bot$ and $\func{clock}()>T_{\mathrm{election}}[\id]+\Delta_{\mathrm{timeout}}$ $\Rightarrow$
\State \quad \textbf{require} $\enabled(\id)$
\State \quad for every $Q\in S_j$: $firstVote[P_j,\id,Q]\gets\timeoutVal_\id$
\State \quad logically emit $\FirstVote(\id,\timeoutVal_\id;S_j)$
\Statex
\State received a valid $\FirstVote(\id,x,\chi)$ from $P_i$ $\Rightarrow$
\State \quad $S_i\gets\EffectiveCommittees(P_i,\id,\chi)$; require $S_i\neq\varnothing$ and $\forall Q\in S_i:firstVote[P_i,\id,Q]=\bot\wedge\neg\finalVoted(P_i,\id,Q)$
\If{$T_{\mathrm{election}}[\id]=\bot$}
    \State $T_{\mathrm{election}}[\id]\gets\func{clock}()$
\EndIf
\State \quad for every $Q\in S_i$: $firstVote[P_i,\id,Q]\gets x$ and add the derived $Q$-local vote to $\mathsf{pool}$
\Statex
\State received a valid $\NotarVote(\id,x,\chi)$ from $P_i$ $\Rightarrow$
\State \quad for every $Q\in\EffectiveCommittees(P_i,\id,\chi)$: add the derived $Q$-local vote to $\mathsf{pool}$
\Statex
\State received a valid $\FinalVote(\id,x,\chi)$ from $P_i$ $\Rightarrow$
\State \quad for every $Q\in\EffectiveCommittees(P_i,\id,\chi)$: add the derived $Q$-local vote to $\mathsf{pool}$
\Statex
\State $\exists x:\admissible(\id,x)$ and $S_j\gets\{Q\in\SignerCommittees(P_j,\id):x\in\manyVotes(Q,\id)\setminus secondLook[P_j,Q,\id]\wedge\firstVoted(P_j,\id,Q)\wedge\neg\finalVoted(P_j,\id,Q)\}$ is nonempty $\Rightarrow$
\State \quad \textbf{require} $\enabled(\id)$
\If{$\id=\Slot(a,s,e)$}
    \State \quad \textbf{require} $\safeToVote_{P_j}(\Slot(a,s,e),x)$
\EndIf
\State \quad for every $Q\in S_j$: $secondLook[P_j,Q,\id]\gets secondLook[P_j,Q,\id]\cup\{x\}$
\State \quad $S'_j\gets\{Q\in S_j:x\notin notarized[P_j,Q,\id]\}$
\If{$S'_j\neq\varnothing$}
    \If{$\id=\Slot(a,s,e)$}
        \State $\votedFor[P_j,\Slot(a,s,e),x]\gets\true$
    \EndIf
    \State broadcast $\NotarVote(\id,x;S'_j)$; for every $Q\in S'_j$: $notarized[P_j,Q,\id]\gets notarized[P_j,Q,\id]\cup\{x\}$
\EndIf
\Statex
\State $S_j\gets\{Q\in\SignerCommittees(P_j,\id):\firstVoted(P_j,\id,Q)\wedge\neg\finalVoted(P_j,\id,Q)\wedge\timeoutAllowed(Q,\id)\wedge\timeoutVal_\id\notin notarized[P_j,Q,\id]\}$ is nonempty $\Rightarrow$
\State \quad \textbf{require} $\enabled(\id)$
\State \quad broadcast $\NotarVote(\id,\timeoutVal_\id;S_j)$; for every $Q\in S_j$: $notarized[P_j,Q,\id]\gets notarized[P_j,Q,\id]\cup\{\timeoutVal_\id\}$
\Statex
\State $\id=\Slot(a,s,e)$ and $(\FirstCert(\id,b)\in\mathsf{pool}\vee\FinalCert(\id,b)\in\mathsf{pool})$ and $b\neq\timeoutVal_\id$ and $\Complete(\id,b)$ $\Rightarrow$
\State \quad $\FinalizeSegment(\id,b)$
\Statex
\State $\id=\Slot(a,s,e)$ and $\exists Q\in\Committees(\Slot(a,s,e)):$ the pool contains valid votes of distinct-signer weight $>F_Q$ in $Q$ for $\Slot(a,s,e)$ $\Rightarrow$
\State \quad $\UpdateVisible(e)$
\end{algorithmic}
\end{breakablealgorithm}

\begin{breakablealgorithm}
\caption{Close-round processing for \(I_r\) at correct replica \(P_j\)}
\begin{algorithmic}[1]
\State use the committee-local receive, vote-counting, fast, final, and timeout-notarization rules for $I_r$
\State $\exists x:$ $S_j\gets\{Q\in\SignerCommittees(P_j,I_r):x\in\manyVotes(Q,I_r)\setminus secondLook[P_j,Q,I_r]\wedge\firstVoted(P_j,I_r,Q)\wedge\neg\finalVoted(P_j,I_r,Q)\}$ is nonempty $\Rightarrow$
\State \quad reconstruct and validate the canonical cut or frontier version for $x$
\State \quad for every $Q\in S_j$: add $x$ to $secondLook[P_j,Q,I_r]$
\State \quad $S'_j\gets\{Q\in S_j:x\notin notarized[P_j,Q,I_r]\}$
\If{$S'_j\neq\varnothing$}
    \State broadcast $\NotarVote(I_r,x;S'_j)$
    \State for every $Q\in S'_j$: $notarized[P_j,Q,I_r]\gets notarized[P_j,Q,I_r]\cup\{x\}$
\EndIf
\Statex
\State $h\neq\timeoutVal_{I_r}$ and $S_j\gets\{Q\in\SignerCommittees(P_j,I_r):\firstVoted(P_j,I_r,Q)\wedge\neg\finalVoted(P_j,I_r,Q)\wedge\func{Support}(P_j,Q,I_r)\subseteq\{h\}\}$ is nonempty and the canonical version for $h$ has been reconstructed and validated $\Rightarrow$
\State \quad \textbf{require} $\enabled(I_r)$
\State \quad for every $Q\in S_j$: $finalVoted[P_j,I_r,Q]\gets\true$
\State \quad logically emit $\FinalVote(I_r,h;S_j)$
\Statex
\State the local vote store derives $\FirstCert(I_r,h)$ for non-timeout $h$ $\Rightarrow$
\State \quad reconstruct and validate the canonical version for $h$
\If{$\closeDecision[I]=\bot$}
    \State $\closeDecision[I]\gets(r,h)$
\ElsIf{$\closeDecision[I].hash\neq h$}
    \State signal a safety violation
\EndIf
\State stop emitting further votes for close instance $I$
\Statex
\State the local vote store derives $\FinalCert(I_r,h)$ for non-timeout $h$ $\Rightarrow$
\State \quad reconstruct and validate the canonical version for $h$
\If{$outcome[I_r]=\pending$}
    \State $outcome[I_r]\gets\final(h)$
\EndIf
\If{$\closeDecision[I]=\bot$}
    \State $\closeDecision[I]\gets(r,h)$
\ElsIf{$\closeDecision[I].hash\neq h$}
    \State signal a safety violation
\EndIf
\State stop emitting further votes for close instance $I$
\Statex
\State the local vote store derives $\NotarCert(I_r,h)$ for non-timeout $h$ $\Rightarrow$
\If{$outcome[I_r]=\pending$}
    \State $outcome[I_r]\gets\converged(h)$
\EndIf
\If{$\closeDecision[I]=\bot$}
    \State enable round $r+1$ using $\func{NextCloseHash}_{P_j}(I,r+1)$
\EndIf
\Statex
\State $\CloseDead_{P_j}(I_r)$ $\Rightarrow$
\If{$outcome[I_r]=\pending$}
    \State $outcome[I_r]\gets\timeoutOutcome$
\EndIf
\If{$\closeDecision[I]=\bot$}
    \State enable round $r+1$ using $\func{NextCloseHash}_{P_j}(I,r+1)$
\EndIf
\end{algorithmic}
\end{breakablealgorithm}

\paragraph{Reconciliation metadata in close-vote gossip.}
A correct replica may advertise retained close-version hashes and available delta bases alongside ordinary
vote gossip.  Such advertisements are unsigned transport metadata unless authenticated by the transport
layer; they are not votes, certificates, certified values, or decisions.  A receiver accepts a
reconstructed version only after applying a delta to a retained canonical base, recomputing the exact
target hash, and running applicable admissibility validation.  Invalid, inapplicable, or unavailable deltas
are ignored, and the receiver asks another replica or obtains a complete canonical version.

\section{Safety}

Assume every weighted committee \(Q\) has Byzantine signer weight at most \(F_Q\), quorum slack
\(P_Q\ge F_Q\), and total signer weight \(W_Q>3F_Q+2P_Q\).  Hence any two quorums of signer weight
at least \(W_Q-F_Q-P_Q\) in the same committee intersect in correct signer weight; fast quorums of
weight at least \(W_Q-P_Q\) only strengthen this.  Each committee in \(\Committees(\id)\) runs the
single-committee discipline independently, and the election outcome is the deterministic merge of its
local terminal results.  Leaves with omitted scope are expanded deterministically into every applicable
committee containing the signer, while explicitly scoped leaves are expanded only into their signed
proper subset; vectorized batches only amortize authentication.  Before a correct replica final-votes
\(x\) in committee \(Q\), its complete support in \((\id,Q)\), including its first vote and every
timeout or ordinary notarization vote counted there, is a subset of \(\{x\}\).  A correct final voter
emits no later first, notarization, or timeout vote in that committee instance, although it may perform
independently valid actions in disjoint committee scopes.

\paragraph{Safety assertions.}
Let \(\mathsf{Correct}\) be the set of correct replicas,
\(\mathsf{Finalized}_j(q)\) the set of non-timeout segment-tip ids that replica \(P_j\) has finalized for
slot election \(q\), and \(\closeDecision_j[I].hash=\bot\) when \(P_j\) has not yet decided logical
close instance \(I\).  The following assertions must hold in every reachable state:
\[
\begin{aligned}
\mathsf{ASSERT}_{\mathsf{slot}}\equiv{}&
\forall P_i,P_j\in\mathsf{Correct},\ \forall q=\Slot(a,s,e),\
\forall b\in\mathsf{Finalized}_i(q),\ \forall b'\in\mathsf{Finalized}_j(q):
b=b',\\
\mathsf{ASSERT}_{\mathsf{cut}}\equiv{}&
\forall P_i,P_j\in\mathsf{Correct},\ \forall e,\
\closeDecision_i[\CloseCut(e)].hash\neq\bot\wedge
\closeDecision_j[\CloseCut(e)].hash\neq\bot\\
&\hspace{8em}\Longrightarrow
\closeDecision_i[\CloseCut(e)].hash=\closeDecision_j[\CloseCut(e)].hash,\\
\mathsf{ASSERT}_{\mathsf{record}}\equiv{}&
\forall P_i,P_j\in\mathsf{Correct},\ \forall e,\
\closeDecision_i[\CloseRecord(e)].hash\neq\bot\wedge
\closeDecision_j[\CloseRecord(e)].hash\neq\bot\\
&\hspace{8em}\Longrightarrow
\closeDecision_i[\CloseRecord(e)].hash=\closeDecision_j[\CloseRecord(e)].hash.
\end{aligned}
\]
The first assertion compares finalization, not merely replica-local terminal observations.  The latter
two compare the canonical value hashes decided for the logical close elections across all rounds;
collision resistance then binds equal hashes to equal canonical cuts or frontier maps.

\paragraph{Same election.}
For a fixed \(\Slot(a,s,e)\), committee-local first-vote uniqueness, \(\manyVotes(Q,\id)\),
\(\timeoutAllowed(Q,\id)\), and the post-final lock satisfy the committee-local kernel defined above.
By the committee-local certificate-incompatibility lemma, a local fast or final certificate for \(b\)
excludes every
local notarization certificate for \(b'\neq b\), including timeout.  A global fast or final certificate
requires the same value in every election committee, so two conflicting local notarization
certificates exclude every possible present or delayed ordinary global final certificate.  The merge
function finalizes a value only when all committee-local results finalize that same value; one timeout
result dominates any non-final compatible result, and different non-timeout values merge to timeout.
Therefore a committee-local result alone has no global finality or durable release effect.  One
matching merged notarization may finish a segment obligation for drain, but remains unresolved until a close-record
hash is certified.  A correct close voter includes its complete segment on a candidate ledger chain only while its
current validated evidence contains exactly one unresolved notarization and no final or release
evidence.  Known conflicting notarizations resolve to release; a hidden conflicting notarization
remains non-final and is released if another segment occupies that starting obligation on the certified
ledger.

\paragraph{Atomic segment finality.}
For one enabled \(q=\Slot(a,s,e)\), every admissible non-timeout value begins at the same certified
base and commits by predecessor hashes to one exact assignment for each account slot in its suffix.  No
election beginning at an internal slot can be enabled while \(q\) is unresolved because
\(s=\NextSlot(a)\).  Same-election certificate incompatibility therefore selects at most one tip, and
that tip selects at most one block for every affected slot.  \(\FinalizeSegment\) validates the complete
suffix in temporary state and publishes every finalized-slot and ledger write in one atomic commit, so
no correct observer can use a partially installed segment as a future base.  After commit,
\(\NextSlot(a)\) advances to one plus the tip sequence.

\paragraph{Release.}
If \(\released_{P_j}(\Slot(a,s,e),b)\), then the certified close state has made \(b\) incapable of
later finalization.  Release follows exactly when the candidate does not occur on the complete account
chains committed by the certified ledger-state root.  Exclusion from the close cut stops drain voting but
does not by itself release a candidate; an ordinary-finalized segment outside the cut remains eligible
and must occur on the close-record frontier.  A timeout or a known pair
of conflicting notarizations supplies a positive witness permitting such omission.  If a hidden
conflicting notarization is disclosed only after a close record including the segment committed by \(b\)
has certified, that segment remains finalized and the hidden candidate is released by certified omission.
Certificate-level incompatibility proves that no segment tip from
a timeout- or conflict-released election can later obtain an ordinary global fast or final certificate,
and uniqueness of the certified close hash prevents a competing close finalization.  The certified-base lock prevents any later segment from extending a released candidate: only the
certified finalized frontier can be the predecessor of the next enabled segment.  Bare timeout or conflict
evidence may finish the segment obligation for drain but is not durable release while the close record remains
unsettled.

\paragraph{Adjacent elections.}
\(\Slot(a,s,e)\) and \(\Slot(a,s,e+1)\) share committee \(\committeeAt(e-2)\).  Every global
non-timeout finishing certificate contains the corresponding local certificate from that shared
committee.  Two conflicting local quorums in the shared committee intersect in a correct replica.
That replica retains its slot-vote history across the epoch boundary, and
\(\safeToVote\) forbids its conflicting later vote unless the earlier obligation has been
certifiably released.  A finishing certificate that still preserves the earlier obligation therefore
contradicts the existence of a conflicting adjacent certificate.

\paragraph{Close visibility, report locks, and cuts.}
A close-cut value becomes available after report weight at least \(W_{Q^\star}-F_{Q^\star}\) in the
deterministic committee \(Q^\star=\ReportCommittee(e)\).  A non-timeout close-cut decision requires a
fast or final result in the unique close committee.  Every correct fresh close-cut voter verifies a
report proof and requires the cut to contain its current \(visible[e]\), \(\ReportVisible(e,\rho)\), and
every election in its own report whose non-finalizability is not proved.  An election already
ordinary-finalized at that voter need not be reported and need not be in the cut.  Historical
\(\func{CutSupported}\) permits convergence on a certified cut even when it omits a low-weight local
report lock, but that omission is not release and the report obligation persists into fresh close-record
validation.

For a global slow fast/final slot certificate in \(Q^\star\), the certificate signer set has weight at
least \(W_{Q^\star}-F_{Q^\star}-P_{Q^\star}\), and its correct signer weight is greater than
\(F_{Q^\star}+P_{Q^\star}\).  At report time, each such correct signer either has already
ordinary-finalized the certified tip and holds an exact finality lock, or has not and reports the
election, acquiring a report lock.  Therefore every correct fresh close-record voter in this combined
set rejects a frontier omitting the certified tip.  The remaining total weight is below the close-record
certificate threshold, so no incompatible close-record hash can obtain persistent support.  A fast slot
certificate only strengthens the bound.  Reaching \(\notarized(b)\) remains sufficient for close drain
when the election is in the cut because the close-record phase will either include its complete segment
or release the obligation by omission.

\paragraph{Safety across close rounds.}
Suppose a valid fast or final certificate for hash \(h\) exists in close round \(I_r\).  Either
certificate decides the logical close instance.  The committee-local certificate-incompatibility lemma
implies that \(I_r\) cannot also have valid timeout evidence or certificate-level conflicting
notarization evidence.  Therefore a correct replica that advances from \(I_r\) without yet learning
the deciding certificate can advance only by carrying \(h\).

In every successor round reached through this live carry, correct replicas first-vote \(h\), even if
another hash is freshly preferred.  Byzantine replicas alone cannot place another value in the
second-look set because their total weight is at most \(F_Q<F_Q+P_Q\), while the second-look threshold
is strictly greater than \(F_Q+P_Q\).  Moreover, once all correct first-vote weight is for \(h\),
\[
\allVotes(Q,I_{r+1})-\maxVotes(Q,I_{r+1})
\le F_Q
\le F_Q+P_Q,
\]
so the ordinary timeout trigger cannot become true solely from Byzantine votes.

Thus no successor round can produce a notarization certificate, timeout certificate, fast certificate,
or final certificate for a value different from \(h\).  By induction, every later close-round fast or
final decision certificate certifies \(h\).

Conversely, if a valid death proof exists for \(I_r\), the committee-local certificate-incompatibility lemma excludes every
fast certificate, and in fact every final certificate, in that round.  Resetting to a fresh hash after
such positive evidence is therefore safe.

The argument does not rely on certificate arrival order.  A valid older-round fast or final certificate
is processed as a decision, even after a replica has entered later rounds.  The instance-wide variable \(\closeDecision[I]\) is
write-once.

\paragraph{Hash-reconstructed close decision.}
For a decided close-cut hash, every correct replica reconstructs the canonical slot set, recomputes the
hash, derives the deciding certificate, and validates the cut through
\(\admissibleHistorical(\CloseCut(e),h_C,C)\).  The certificate attests, under the committee fault
bound, that correct weight freshly validated the report quorum and every start obligation before voting;
those local witness choices need not be serialized or reconstructed after certification.  For a decided
close-record hash, every correct replica reconstructs the canonical frontier map, recomputes the hash
using the preceding close-record hash, and validates the historically admissible decided cut,
\(\RecordSupported\) branch, slot evidence for both cut elections and ordinary-finalized epoch-\(e\)
elections outside the cut, complete account chains, ledger transitions, ancestry, and
duplicate-account-slot constraints.  Evidence learned after a vote may change fresh preference but
cannot change the canonical preimage of a decided hash.

\paragraph{Soundness of reconciliation.}
A reconciliation delta has no consensus authority.  It is accepted only relative to a locally retained
canonical base, and only when applying the delta produces a canonical set or frontier map whose hash is
exactly the target and whose applicable fresh or historically supported admissibility checks succeed.  A malformed delta therefore either
fails application, produces the wrong hash, or fails semantic validation.  Subject to collision
resistance of \(\hash\), it cannot make a correct replica accept contents different from the contents
committed by the target hash.  Applying a delta never deletes underlying blocks, reports, or votes and
does not alter the replica's current report store, vote store, or derived visibility; it only reconstructs
an immutable historical version indexed by hash.

\paragraph{Disjoint elections.}
For \(e'\ge e+2\), elections need not share a committee.  Every valid vote in epoch \(e'\), however,
is interpreted under the implicit governing context
\(\governingHash(e')=\closeHash[e'-2]\), and the voter must validate the corresponding cumulative
\(\closeState[e'-2]\) and its certified account frontiers.  That state contains only finalized or
released slot outcomes.  If it finalizes block \(b\) as the frontier of account \(a\),
\(\safeFromClose\) permits only a segment beginning at the next slot whose tip descends from \(b\) and whose complete suffix validates from that certified base.  If the state instead contains \(\ReleasedState(\pi)\), the
certified close decision releases the earlier same-starting-slot obligation, while the certified-base lock forbids a released block from becoming the base or an internal block of a later certified segment.  Since every unresolved
notarization is either finalized or released when its close record certifies, no unresolved obligation
is carried across disjoint committee generations.

A correct replica may vote in epoch \(e'\) only after it has processed the certified close state of
\(e'-2\), including durable \(\closeHash[e'-2]\), \(\ledgerRoot[e'-2]\), \(\closeCut[\cdot]\),
\(\frontier[\cdot,\cdot]\), release evidence, and finalized-slot state.  The close state therefore
participates directly in vote validity rather than only indirectly in committee derivation.

\paragraph{Theorem.}
\(\mathsf{ASSERT}_{\mathsf{slot}}\), \(\mathsf{ASSERT}_{\mathsf{cut}}\), and
\(\mathsf{ASSERT}_{\mathsf{record}}\) hold.  In particular, no two correct replicas finalize
conflicting block ids for the same account slot.  If conflicting slot finality came from the same
election, same-election safety contradicts it; from adjacent elections, adjacent safety contradicts it;
and from elections at distance at least two epochs, disjoint-election safety contradicts it.  Close-cut
and close-record agreement follow from certificate incompatibility, live carry across rounds, and the
write-once logical close decision.


\section{Liveness}

The liveness statement distinguishes a terminal slot-election observation from a durable slot
decision.  A notarization may finish the slot election for drain, but the slot is settled only when an
ordinary fast or final certificate finalizes it or when the certified ledger includes its candidate or
releases the election by omission.  Define
\begin{center}
\begin{adjustbox}{max width=\columnwidth}
\(\displaystyle
\begin{aligned}
\func{settled}(\Slot(a,s,e))\iff{}&
\finalizedSlot[a,s]\neq\bot\\
&\vee\bigl(state[e]=\mathsf{closed}\wedge\closeCut[e]\neq\bot\\
&\qquad\wedge(\Slot(a,s,e)\notin\closeCut[e]\vee
\closeState[e](\Slot(a,s,e))\in\{\FinalizedState(b,\pi),\ReleasedState(\pi)\})\bigr).
\end{aligned}
\)
\end{adjustbox}
\end{center}

\paragraph{Liveness assertion.}
With \(\Diamond\) denoting ``eventually'' over an admissible execution satisfying the liveness
assumptions below, epoch closure is asserted replica-by-replica:
\[
\mathsf{ASSERT}_{\mathsf{epoch\mbox{-}closure}}\equiv
\forall e\ge0,\ \forall P_j\in\mathsf{Correct}:\quad
\Diamond\bigl(state_j[e]=\mathsf{closed}\bigr).
\]
This is a temporal progress assertion, not a finite-prefix invariant: it is evaluated only for
executions with the stated eventual-synchrony, responsiveness, gossip, finite-evidence, retention, and
reconstruction assumptions.
For a logical close instance \(I\in\{\CloseCut(e),\CloseRecord(e)\}\), settled means
\(\closeFinished(I)\); a round-local \(\converged(h)\) observation is not settlement.  Thus
conditional election liveness for RAI is
\begin{center}
\begin{adjustbox}{max width=\columnwidth}
\(\displaystyle
\started[\id]\Rightarrow \Diamond\func{settled}(\id),
\)
\end{adjustbox}
\end{center}
where slot elections use the definition above and logical close instances use
\(\func{settled}(I)\iff\closeFinished(I)\).

\paragraph{Close-liveness assumptions.}
Assume:

\begin{enumerate}
\item after some time the network is synchronous;
\item every correct committee member remains responsive long enough to perform every committee-local
      vote action once that action is enabled in that committee;
\item whenever the committee-local termination argument requires a correct signer to issue a first,
      second-look, timeout-notarization, or final vote in one applicable committee, the signer eventually
      emits a vote whose effective scope contains that committee; the scope is omitted when the same
      action applies to every applicable committee and is explicit when only a proper subset can act;
\item every valid block, report, and vote learned by one correct online replica is eventually made
      available to every correct online replica through duplicate-suppressing reliable gossip;
\item a correct replica that votes for a close hash retains its canonical cut or frontier preimage until
      the close instance decides;
\item decided versions, or valid delta paths from retained checkpoints, remain retrievable from correct
      archival replicas;
\item reconstruction and validation are objective, so every correct replica can eventually reconstruct
      and validate a carried hash;
\item the block, report, and vote universe relevant to elections in the frozen certified cut and to
      ordinary-finalized epoch-\(e\) segments outside that cut is finite, or eventually stops changing the
      canonical fresh close record;
\item every fresh report lock for an election omitted from the decided cut is eventually satisfied by a
      support-compatible included segment, discharged by positive timeout/conflict evidence, or superseded
      by persistent close-record support whose quorum-intersection proof excludes hidden ordinary finality;
      and
\item fresh construction and canonical serialization are deterministic.
\end{enumerate}

The parameter \(P_Q\) provides quorum-weight slack but does not independently provide crash tolerance for
this liveness theorem.

\paragraph{Common validation of a carry.}
A non-timeout notarization certificate is locally derived from enough authenticated vote leaves that at least one
correct voter reconstructed and validated the corresponding close version.  Under version
retention, reconciliation availability, and complete atom gossip, every correct replica eventually
reconstructs that same canonical version.  For a carried cut hash, the underlying same-hash committee-local certificates make the historical
admissibility branch permanent; for a carried frontier hash, \(\RecordSupported\) makes the historical
branch permanent.  The lock-intersection result ensures that a supported frontier cannot omit hidden
ordinary finality.  Every correct replica can therefore validate and first-vote the carried hash even if
its fresh report/finality preference differs.

\paragraph{Relative reconstruction of a carry.}
A correct replica obtains the canonical preimage of a live carried hash by applying an add/remove cut
delta or sparse frontier replacements to a retained base, or by retrieving a complete version when no
usable base is available.  It recomputes the target hash and evaluates fresh admissibility when no
support certificate exists, or historical admissibility for a supported cut or frontier version.  Failure, delay, or corruption of one reconciliation response cannot prevent
progress because the replica may ask another correct holder or obtain the complete canonical version.

\paragraph{Eventual fresh-choice convergence.}
Let \(E_j(t)\) be the relevant valid evidence known to correct replica \(P_j\).  Complete gossip and
the finite-evidence assumption imply that the correct replicas eventually have the same effective
evidence for fresh construction.  Canonical construction therefore gives
\[
\exists T,h^\ast\quad
\forall t\ge T,\ \forall P_j\text{ correct}:
\FreshCloseHash_{P_j}(I)=h^\ast,
\]
whenever no live carry constrains the next first vote.

Correct replicas need not have identical complete sets of admissible hashes.  They need only
eventually validate every live carry and eventually select the same fresh hash in an unlocked round.

For state synchronization and bootstrapping, a joining replica may use the next live close as the latest
target but obtains the durable blocks, reports, and vote batches and reconstructs the decided cut and frontier
versions for every epoch from genesis onward.  For each epoch it validates all blocks and transitions,
locally derives the deciding close-vote quorums, recomputes the close-cut hash, and recomputes the
close-record hash from the preceding close-record hash and the canonical frontier map.  The terminal
validated close commits transitively to the certified block chains and yields the live committee by
replay.  The archive supplying atoms or reconciliation deltas is not trusted.

A replica starts closing an open epoch only after the previous epoch and its close state are
certified; when it starts closing that epoch, it immediately opens the successor epoch.  Every
responsive participant in epoch \(e\) eventually obtains \(\closeState[e-2]\) before voting.  When
a decided close cut \(C\) is installed, elections in \(C\) are resumed for the drain phase until they
finish.  Elections outside \(C\) are not resumed; their ordinary-finalized segments remain eligible for
the epoch-\(e\) close record, and all other candidates are settled only by certified close-record
omission.  The complete-gossip, finite-evidence, and report-lock-resolution assumptions above govern
fresh close-record construction.

\paragraph{Election-core termination.}
Fix an enabled committee instance \((\id,Q)\), and write
\(W=W_Q\), \(F=F_Q\), \(P=P_Q\), \(g=F+P\), and \(q=W-F-P\).
Assume a sufficiently long synchronous period in which every correct member of \(Q\) remains
online long enough to send its first vote, receive every correct first vote, and perform every enabled
second-look or timeout action.  Byzantine members have total weight at most \(F\), and
\(W>3F+2P\).  Then the committee instance eventually derives one local terminal result
\[
L_Q(\id)\in\{\timeoutOutcome,\notarized(x),\fast(x),\final(x)\}.
\]

Let \(F'\le F\) be the actual Byzantine weight.  After all correct first votes have been delivered,
exactly one of the following cases holds.

First, suppose some non-timeout value \(x\) has valid first-vote signer weight strictly greater than
\(g=F+P\).  At most \(F\) of that weight is Byzantine, so the correct first-vote weight supporting
\(x\) is strictly greater than \(P\).  A correct replica emits a non-timeout first vote only after
obtaining the block data or reconstructing the close version and validating \(\admissible(\id,x)\).
Under complete atom gossip, version retention, and reconciliation availability, every correct member of
\(Q\) eventually obtains the data required to validate \(x\).  For a slot value, \(\ExtendsTree\),
\(\safeFromClose\), and the enabled retry guards depend only on the validated block and ancestry,
monotone finalized state, and certified governing close state.  For a close value, fresh admissibility may differ because of signer-local report or finality locks, but
persistent same-hash support creates a stable historical branch.  If a popular unsupported hash is
compatible with the converged locks, every correct member eventually performs the committee-local second
look and issues a notarization vote for \(x\).  If it conflicts with an actual ordinary-finality lock,
correct refusing weight is greater than \(F+P\), so the hash cannot obtain support and the opposing vote
weight enables timeout or convergence on the including hash.  Low-weight report locks are resolved under
the stated close-liveness assumption or safely superseded only after persistent support.  Correct signer weight is at least \(W-F'\ge W-F\ge W-F-P=q\), so
\(Q\) derives a local notarization certificate for \(x\).

Second, suppose no non-timeout value has first-vote signer weight strictly greater than \(g\).  Once all
correct first votes are present,
\[
\allVotes(Q,\id)\ge W-F'
\qquad\text{and}\qquad
\maxVotes(Q,\id)\le F+P.
\]
Therefore
\[
\allVotes(Q,\id)-\maxVotes(Q,\id)
\ge (W-F')-(F+P)
\ge W-2F-P
>F+P.
\]
Thus \(\timeoutReady(Q,\id)\) becomes true.  Every correct member that has not final-voted in this
committee instance eventually issues the enabled timeout notarization vote; if a correct member has
already final-voted, the corresponding local final result is already obtainable.  Hence in either case
\(Q\) eventually derives a local final result or a local timeout certificate.  This is a
stake-weighted projection of the two cases in Kudzu's Liveness~II lemma
\cite[Lemma~5.10]{shoup2025kudzu}; the additional RAI assumptions above supply stable admissibility and
close-version reconstruction.

Every signed vote counts in exactly the committee instances in its effective scope, using the signer's
frozen weight in each such committee.  Because a correct signer may emit independently valid actions for
disjoint committee scopes, an action required by one committee is not suppressed merely because another
applicable committee requires a different second-look, timeout, or final action.  The committee-local
termination argument therefore applies independently to every \(Q\in\Committees(\id)\) without a
joint-action compatibility assumption.  If different committees terminate on different non-timeout
values, their certificates merge to timeout and conflict evidence rather than global finality.

By the committee-composition lemma, after every committee has a local terminal result, the
deterministic merge immediately finishes the current election round.  For a slot election, notarization
in every committee yields \(\notarized(x)\).  For a close election, matching notarizations yield
\(\converged(x)\), which starts a carried-value successor round rather than certifying the close.
Timeout in one committee and a non-timeout result in another yields timeout; matching final results for
the same value yield \(\final(x)\).  For a close election, that result is a decision.  Different
non-timeout values, including conflicting committee-local final values, merge to timeout and their
signed terminal certificates form conflict death evidence that unlocks a later round.

The committee-local termination argument adapts only Kudzu's two-case vote-counting structure; RAI's
values are derived implicitly as above, and first votes are the first commitment to the derived value.
For close elections, a matching non-timeout notarization certificate creates a live carry, a valid fast
or final certificate decides the logical instance, and timeout or conflict death evidence unlocks a
later round.

\paragraph{Safe close attempts after the reporting-committee quorum.}
Before a non-timeout fresh close-cut vote, a correct replica locally selects a report proof containing
valid eligible signed reports of combined signer weight at least \(W_{Q^\star}-F_{Q^\star}\) in
\(Q^\star=\ReportCommittee(e)\).  It requires its canonical cut version to include
\(\ReportVisible(e,\rho)\), its current local visibility, and every election under its own fresh report
lock.  The proof is not transmitted with the close vote; every receiver derives a suitable proof from
its gossiped report store.

Let \(S^\star\) be the signer set of any global slow fast/final certificate for
\(q=\Slot(a,s,e)\) in \(Q^\star\).  Its correct signer weight is strictly greater than
\[
(W_{Q^\star}-F_{Q^\star}-P_{Q^\star})-F_{Q^\star}
= W_{Q^\star}-2F_{Q^\star}-P_{Q^\star}
>F_{Q^\star}+P_{Q^\star}.
\]
At the instant reports are signed, every correct member of \(S^\star\) is in exactly one of two safety
classes: it has already reconstructed and applied ordinary finality for the certified tip, or it has not
and therefore reports the still-finalizable election because its non-timeout vote remains in support.
The first class freshly rejects every close-record frontier omitting the exact certified tip; the second
class freshly rejects omission under its report lock and, after learning finality, adopts the same exact
lock.  Thus more than \(F_{Q^\star}+P_{Q^\star}\) correct weight rejects an incompatible fresh
close-record hash.  Such a hash cannot reach the local notarization, fast, or final threshold and hence
cannot obtain \(\RecordSupported\).  This proves that a historically supported close record may safely
override low-weight report locks: if it omits an election, no hidden ordinary fast/final certificate for
that election can coexist with the support.

Consequently, any close-cut hash stored in the instance-wide close decision safely freezes the drain cut
\begin{center}
\begin{adjustbox}{max width=\columnwidth}
\(\displaystyle
C=\cutVersion[e,h_C].
\)
\end{adjustbox}
\end{center}
Every unresolved election forced by report or vote visibility is in \(C\); an ordinary-finalized
election may be outside \(C\) and is incorporated directly into the close-record frontier.  An election
outside \(C\) is not released by cut exclusion.  It is finalized if its ordinary-certified segment is
on the decided frontier and otherwise is released only by certified close-record omission.

\paragraph{Visibility and fresh preference across close attempts.}
Vote-derived visibility is monotone, while report-derived visibility is recomputed from the unique
eligible report of each non-equivocating signer.  Because correct replicas gossip all valid reports and
votes, evidence seen by one correct online replica eventually becomes available to every other correct
online replica, and every correct replica eventually excludes the same equivocating report signers.
In an unlocked round, this may change \(\FreshCloseHash_{P_j}(\CloseCut(e))\) or
\(\FreshCloseHash_{P_j}(\CloseRecord(e))\).  It does not delete a retained historical candidate.
Unsupported candidates are freshly revalidated under the current eligible-report and lock state before
a new vote; a cut or record with persistent certificate support is instead validated by the historical
branch and never loses validity.  A live carry overrides fresh preference whenever its carried version
remains admissible.  Under the
finite-evidence and canonical-construction assumptions,
fresh preferences eventually converge whenever no carry constrains the next first vote.

\paragraph{Close-round termination.}
Consider a close round \(I_r\).

If a non-timeout hash \(h\) obtains first-vote signer weight strictly greater than \(F_Q+P_Q\) in an
applicable committee and is compatible with the converged fresh report/finality locks, atom gossip and
version reconciliation allow every correct committee member to perform the second look for \(h\).  If
an unsupported hash conflicts with fresh locks, the report-lock-resolution assumption eventually either
makes it compatible, supplies positive death evidence, or yields a different including hash; if it
obtains persistent same-hash support first, the historical branch makes it common.  An actually
ordinary-finalized segment cannot be lost in this transition because its correct lock weight is greater
than \(F_Q+P_Q\) and prevents incompatible support.  Thus the popular-candidate case cannot deadlock
merely because another hash has become freshly preferred.

If no non-timeout hash obtains first-vote signer weight strictly greater than \(F_Q+P_Q\), then after
all correct first votes are delivered,
\[
\allVotes(Q,I_r)-\maxVotes(Q,I_r)>F_Q+P_Q,
\]
so the committee-local \(\timeoutReady\) rule becomes enabled.

Therefore every unlocked round eventually:
\begin{itemize}
\item produces a deciding fast or final certificate;
\item produces a live non-timeout carry; or
\item produces a death proof.
\end{itemize}

If it produces a live carry \(h\), every correct replica eventually validates \(h\) and first-votes it
in the successor round.  Since \(P_Q\ge F_Q\),
\[
W_Q-F_Q\ge W_Q-P_Q,
\]
so the correct replicas alone reach the fast threshold and decide \(h\).

If it produces a death proof, fresh voting resumes.  Several unlocked rounds may occur before evidence
convergence, but each terminates.  Eventually an unlocked round begins after all correct replicas
compute the same fresh hash \(h^\ast\); the correct first-vote weight then reaches the fast threshold
and decides \(h^\ast\).

\paragraph{Draining the decided cut.}
Once \(\closeDecision[\CloseCut(e)]=(r,h_C)\), replicas reconstruct and install
\begin{center}
\begin{adjustbox}{max width=\columnwidth}
\(\displaystyle
C=\cutVersion[e,h_C],\qquad \closeCut[e]=C.
\)
\end{adjustbox}
\end{center}
For every \(q=\Slot(a,s,e)\in C\), the historically admissible certified cut preserves the slot's start
obligation, and installation initializes \(T_{\mathrm{election}}[q]\) if necessary.  Thus either the slot election has already finished, or it
is resumed with a running timeout during the close-drain phase.  By election-core termination, each
such included slot election eventually
finishes, possibly with \(\notarized(b)\) or \(\timeoutOutcome\).  For every \(\Slot(a,s,e)\notin C\), the election is not resumed after the certified cut.  Its
ordinary-finalized segment remains eligible for the close record; otherwise the election is settled by
certified close-record omission when epoch \(e\) closes.

When epoch \(e\) enters closing, \(\StartClosing(e)\) immediately opens epoch \(e+1\), subject to
the invariant that epoch \(e-1\) is already closed.  The decided cut then fixes
\(\mathsf{CloseInput}_e=(\closeHash[e-1],\frontier[e-1,\cdot],C)\).  After all elections in \(C\) have
finished, each replica applies the deterministic resolver only to those elections, adds every validated
ordinary-finalized epoch-\(e\) segment outside \(C\), and performs ancestry closure by close-local replay
from \(\closeState[e-1]\), deriving a candidate frontier for every account.  Ordinary finalizations from
epoch \(e+1\) do not change the epoch-\(e\) candidate frontier.  A global fast or final certificate requires the complete segment committed by its tip to occur on the
candidate ledger from a certified base; locally derived timeout or conflicting-notarization evidence
omits the segment obligation; and exactly one currently known non-timeout notarization, with no known
final or release evidence, may place its complete segment on the candidate account chain.  There is no tie-breaking when a conflict is known.

The evidence known to a correct replica grows monotonically.  If no live carry constrains the next
round, complete atom gossip and finite evidence eventually give all correct online replicas the same
effective evidence for fresh construction, so they derive the same canonical frontier map \(\Phi_e\)
and the same close-record hash
\[
\hash(\mathsf{RAI/CloseRecord}\parallel e\parallel\closeHash[e-1]
\parallel\operatorname{canon}(\Phi_e)).
\]
If a close-record fast or final decision is derived in any round, the instance-wide decision is authoritative:
every block on the decided account chains finalizes and every candidate omitted from those chains is
released, whether its election was inside or outside the cut, even if conflicting notarization or timeout
votes are learned later.  Once
\(\closeDecision[\CloseRecord(e)]=(r,h_F)\), the epoch loop reconstructs the decided frontier version,
validates and installs \(\closeState[e]\) and \(\frontier[e,\cdot]\), sets
\(\ledgerRoot[e]\gets\hash(\operatorname{canon}(\Phi_e))\) and \(\closeHash[e]\gets h_F\), derives
\(\committee[e]\) from the installed balances and representatives, and calls \(\FinishClosing(e)\).  Thus epoch \(e+1\) opens when epoch \(e\)
starts closing, and epoch \(e\) eventually becomes closed after its certified cut has reached terminal
outcomes and its compact close record has been decided.


\paragraph{Theorem.}
Under the assumptions above, every started slot election is eventually settled and every started
logical close instance eventually decides.  A close round either decides, creates a live carry, or
produces a death proof.  A live carry is validated and fast-decided in its successor round; a death
proof unlocks canonical fresh voting, which eventually converges and reaches the fast threshold.
Valid fast or final certificates are processed as decisions in every round, including after a replica
has entered a later round.  The decided close cut is
drained, the hash-reconstructed decided close record installs the certified account frontiers and all
derived finalized and released states, and the successor epochs continue to open and close during
periods of synchrony.  Consequently,
\(\mathsf{ASSERT}_{\mathsf{epoch\mbox{-}closure}}\) holds.

\balance
\begin{thebibliography}{1}

\bibitem{shoup2025kudzu}
Victor Shoup, Jakub Sliwinski, and Yann Vonlanthen.
\newblock 2025.
\newblock Kudzu: Fast and Simple High-Throughput BFT.
\newblock \emph{arXiv preprint arXiv:2505.08771v2} (29 September 2025).

\end{thebibliography}
\end{document}
